\chapter*{Preface}
\addcontentsline{toc}{chapter}{Preface}

This book develops analysis by constructing things: rational enclosures for
numbers, polygonal areas and solid volumes, trigonometric values, rectangle
integrals, secant computations of derivatives, series and local models,
finite evolutions, and transforms with proved endpoint errors.  Every
infinite process comes with a proof that its finite computations agree
and become arbitrarily accurate.

The emphasis is on useful examples and reusable constructions, not on the
largest possible classes of set-theoretic functions.  Rational inputs do
not force rational outputs.  Continuity data permits evaluation on computed
inputs without making a completed real line the foundation of the subject.
An identity means equivalence of two computations, often obtained in very
different ways.

After the interval foundations, geometry comes before trigonometry.  Inner
and outer polygons compute the circle; finite cylinder stacks and conical
frusta establish the sphere--cylinder volume and surface comparisons.
These geometric exhaustions require neither angle functions nor a general
integral.  A separate chapter then uses sector area to construct angles,
sine, and cosine.  The subsequent route is
\[
 \begin{gathered}
 \text{independent rectangle integrals and convex secant derivatives}\\
 \longrightarrow\ \text{the FTC, series, and elementary computations}\\
 \longrightarrow\ \text{local models, complex paths, and parameter integrals}\\
 \longrightarrow\ \text{Fourier computations, impulses, and fundamental matrices}\\
 \longrightarrow\ \text{special functions, transforms, and equivalent representations}.
 \end{gathered}
\]

There is a reason for this order.  A function can be continuous and
differentiable at every rational point while hiding a jump at an irrational
cut.  The fundamental theorem must not mistake its pointwise derivative
for all its increments.  Convex secant inequalities supply the missing
interval control and construct the derivative itself.  Explicit jumps
can subsequently be retained as delta inputs, rather than discarded.

A local series brings higher derivatives and remainder bounds.  A
nonlinear inverse brings a contraction estimate.  An improper integral
brings endpoint tails, and its parameter derivatives bring an integrable
majorant.  These different descriptions are not forced into one global
function class.  The explicit smooth cutoff shows why analyticity is not
a universal assumption; Lambert's function and gamma show why a linear
ODE is not a universal representation either.

The fundamental matrix unifies accumulation and propagation.  An ordinary
forcing term is integrated against it; an impulse is a finite state update
propagated by the same matrix.  Bessel and hypergeometric recurrences supply
local data which can be transported along certified paths.  The sinc and
Gaussian transforms, gamma recurrence, and Euler--Maclaurin construction
of zeta provide other ways to enlarge the domain of useful computations.
An elliptic integral closes the main route by connecting geometry,
quadrature, a series, and a differential equation.

The mathematical prerequisites for each operation are stated where they
are used and justified by finite estimates.  Theorems about unrestricted
function or distribution spaces do not supply missing computations.
Conversely, a useful finite construction is not held back because a more
general theory could also contain it.  The aim is that working with a
specific elementary or special function has a clear route through the
same foundations, with new estimates added where that function needs them.

\clearpage
\chapter{Computable Numbers and Functions}\label{ch:foundations}

A number will be a computation, not an already available point of a
completed line.  The computation gives rational bounds, and a proof explains
why they agree and how they become arbitrarily close.  Different computations
of the same number remain worth keeping: an elementary construction can
explain a quantity which a different construction computes more quickly.

The finite arithmetic used below is ordinary integer and rational arithmetic.
Fractions are reduced by the Euclidean algorithm; finite sums, products,
polynomials, induction, and the triangle inequality require no limiting
objects.  A rational input is a numerator and a positive denominator, not an
oracle for a point on a continuum.

\section{Interval arithmetic and real numbers}

\begin{definition}[Rational interval]\label{def:qinterval}
A \emph{rational interval} is a pair
\[
 [a,b]\in\Q\times\Q,\qquad a\le b.
\]
It is a finite object, not the set of rational points between its endpoints.
Its width is $b-a$.  Two intervals $[a,b]$ and $[c,d]$ overlap when
$a\le d$ and $c\le b$.
\end{definition}

\begin{definition}[A computed number]\label{def:raw-real}
A number computation is a terminating algorithm
\[
 n\longmapsto X_n=[a_n,b_n]
\]
whose intervals satisfy
\[
 n\le m\ \Longrightarrow\ a_n\le a_m\le b_m\le b_n,
\]
and whose widths become smaller than every positive rational tolerance.
We also call such an algorithm a raw real.  No infinitely stored table of
rational endpoints is allowed.
\end{definition}

A request for precision $\varepsilon>0$ is answered by computing successive
stages until $b_n-a_n<\varepsilon$.  A proof that the widths shrink proves
that this search terminates.  An explicit rate is often useful, but need not
be part of the algorithm's input.

\begin{definition}[Equivalence]\label{def:raw-equiv}
Two valid computations are equivalent when their intervals overlap at every
stage.  We write $X\sim Y$, and use the usual equality sign in mathematical
identities between their values.
\end{definition}

Same-stage overlap implies overlap at any two stages, by passing to their
maximum.  Equivalence is transitive: if two outer computations had disjoint
intervals separated by a positive rational gap, a sufficiently narrow interval
of a middle computation could not overlap both.  This proves transitivity
using only finitely many rational comparisons.

Equality is a theorem about computations, not a proposed decision algorithm.
To show two values differ, disjoint finite enclosures suffice.  To show a
value is positive, it suffices to find an enclosure with positive lower
endpoint.  Such a witness is called a separation from zero.

\begin{definition}[Interval arithmetic]\label{def:raw-real-arithmetic}
Addition and negation are computed by
\[
 [a,b]+[c,d]=[a+c,b+d],\qquad -[a,b]=[-b,-a].
\]
For multiplication take the least rational interval containing
$ac,ad,bc,bd$.  If an interval $[a,b]$ avoids zero, its reciprocal enclosure
is $[1/b,1/a]$.
\end{definition}

These operations preserve inclusion.  If the factors have absolute bounds
$B_X,B_Y$, the product width is at most
\[
 B_X\operatorname{width}(Y)+B_Y\operatorname{width}(X).
\]
The bounds can be read from any initial enclosures.  If $|X|\ge c>0$ is
certified, reciprocal widths are at most $c^{-2}\operatorname{width}(X)$
once the intervals have that separation.  These estimates prove validity
and compatibility with equivalence.  In particular, division is available
when a nonzero denominator has been certified, not by deciding an arbitrary
computed value's equality with zero.

An inequality such as $|X-Y|\le r$ has an interval meaning: arbitrarily
accurate enclosures certify the bound $r+\eta$ for each rational $\eta>0$.
An assertion of containment in a closed rational interval can equivalently
be phrased as overlap with that interval at every stage.  None of this
requires choosing a point in an infinite intersection.

\begin{proposition}[Finite-prefix stabilization]\label{prop:raw-prefix-stabilization}
Suppose rational numbers $q_n$ and positive rational errors $e_n\to0$
are computed, and
\[
 |q_m-q_n|\le e_n\qquad(m\ge n).
\]
Then
\[
 X_n=\bigcap_{k=0}^n[q_k-e_k,q_k+e_k]
\]
is a valid number computation.
\end{proposition}
\begin{proof}
The finite intersection contains $q_n$.  It is nested, and its width is at
most $2e_n$.  All its endpoints are finite maxima and minima of rationals.
\end{proof}

More generally, a sequence of pairwise overlapping rational intervals with
shrinking widths can be stabilized in this way: pairwise overlap of a
\emph{finite} family of intervals implies a nonempty common intersection.
We shall repeatedly prove compatibility by common refinements or by bounds
on every finite future correction.  An independently existing limit is
never needed as an anchor.

\begin{definition}[Stage schedules]\label{def:stage-schedule}
A stage schedule is an increasing, unbounded map $\sigma:\N\to\N$.
Replacing $X_n$ by $X_{\sigma(n)}$ gives an equivalent computation.
\end{definition}

Thus two algorithms can be compared at different speeds.  A useful named
number can carry several computations connected by proved equivalences;
transitivity supplies comparisons along a finite chain.  There is no reason
to discard a slow geometric construction after discovering a faster series.

\section{Examples of real numbers}

\subsection{Square roots and finite geometry}
\label{rem:algebraic-numbers-sqrt-two}
For a rational $q\ge0$, start with rational $0\le a\le b$ satisfying
$a^2\le q\le b^2$.  Bisect and retain the half whose squared endpoints
still bracket $q$.  If the midpoint square equals $q$, return the exact
constant interval thereafter.  At stage $n$ the width is at most
$(b-a)2^{-n}$, and finite interval multiplication proves that the resulting
computation squares to $q$.

For $\sqrt2$, the rational secant--tangent refinement gives another
construction:
\begin{verbatim}
def sqrt2(n):
    a = 1
    b = 2
    for _ in range(n):
        a = (2 + a*b) / (a + b)
        b = (b + 2/b) / 2
    return [a, b]
\end{verbatim}
All divisions here are exact rational divisions.  For one step, writing
$a',b'$ for the new endpoints,
\[
 (a')^2-2=\frac{(a^2-2)(b^2-2)}{(a+b)^2},\qquad
 (b')^2-2=\frac{(b^2-2)^2}{4b^2}.
\]
Also $a\le a'\le b'\le b$, and
\[
 b'-a'=\frac{(b-a)(b^2-2)}{2b(a+b)}\le\frac{b-a}{2}.
\]
These identities prove validity.  The bisection and secant--tangent
computations are equivalent because both retain the unique nonnegative
solution, with uniqueness following from $(u-v)(u+v)=u^2-v^2$ and the
positive lower bounds in this example.

\begin{figure}[htbp]\centering
\sqrtSecantTangentAnimation
\caption{Secant--tangent rational enclosures for $\sqrt2$.}
\label{fig:sqrt-secant-tangent}
\end{figure}

There is no rational whose square is $2$: in a reduced fraction $p/q$, the
identity $p^2=2q^2$ would make both $p$ and $q$ even.  Thus this is a
nonrational quantity constructed entirely by rational arithmetic.

Square roots also compute lengths of rational vectors.  The finite identity
\[
 (a^2+b^2)(c^2+d^2)-(ac+bd)^2=(ad-bc)^2
\]
proves the coordinate Cauchy--Schwarz inequality and hence the broken-line
inequality.  Shoelace determinants compute polygonal areas exactly.  These
are the geometric tools needed for the circle, rather than a separate
axiomatization of Euclidean space.

\subsection{Other computations of numbers}
\label{rem:sources-of-raw-reals}
The geometric sum illustrates a finite identity with a tail:
\[
 (1-r)\sum_{k=0}^{N}r^k=1-r^{N+1}.
\]
For rational $|r|<1$, the error $|r|^{N+1}/(1-|r|)$ gives a number
computation.  Bounded monotonicity by itself does not supply a shrinking
upper--lower gap; an error estimate or a terminating precision search is
still necessary.

Later we shall construct, and prove equivalent, the computations suggested
by
\[
 \frac\pi4=1-\frac13+\frac15-\cdots
   =\int_0^1\frac{dt}{1+t^2},\qquad
 \sum_{n=1}^{\infty}\frac1{n^2}=\frac{\pi^2}{6}.
\]
At this point these are destinations, not definitions silently imported from
ordinary calculus.  The circle will define $\pi$ geometrically; rectangles,
series, and Fourier sums will then supply different computations of it.

\begin{remark}[Calculus as equivalence]\label{rem:calculus-as-equivalence}
A great deal of calculus is the art of proving that independently constructed
numbers are equivalent.  The fundamental formula
\[
 \int_a^b f'(x)\,dx=f(b)-f(a)
\]
turns an accumulation computation into an endpoint computation.  Its
hypotheses must explain why local derivative information controls a whole
rational interval; pointwise differentiation on the rationals alone does not.
\end{remark}

\section{Complex numbers}

Let $\Q(i)=\{r+is:r,s\in\Q\}$.  A rational complex box is a pair of
rational coordinate intervals.  A complex computation gives nested boxes
whose two widths shrink to zero; equivalence is coordinatewise overlap.
Addition and multiplication are rational box operations, using
\[
 (r+is)(u+iv)=(ru-sv)+i(rv+su).
\]
For a denominator separated from zero use
$z^{-1}=\overline z/( (\Re z)^2+(\Im z)^2)$ and the real reciprocal
construction.  Finite rational coordinate bounds control all errors.

The same convention applies to vectors and matrices: finitely many number
computations, with finite sums and products.  No additional completeness
principle appears when the number of coordinates increases.

\section{Algorithms with a variable}
\label{rem:parameterized-constructions}

\begin{definition}[Function computation]\label{def:real-fun-raw}
A real-valued function computation consists of a declared domain of rational
inputs and an algorithm
\[
 (x,n)\longmapsto [a_f(x,n),b_f(x,n)]
\]
which is a valid number computation at each permitted input $x$.
\end{definition}

The outputs need not be rational.  In particular, $\sin(\pi x)$ will have
rational inputs and computed outputs.  An algebraic, dyadic, or rational
geometric chart may be a more natural first domain than all rationals;
passing between charts requires an explicit comparison of computations.

\begin{definition}[Interval continuity]\label{def:interval-regular}
On a rational interval $[a,b]$, a continuity certificate supplies, for every
positive rational $\varepsilon$, a positive rational $\mu(\varepsilon)$
such that
\[
 |x-y|\le\mu(\varepsilon)\quad\Longrightarrow\quad
 |f(x)-f(y)|\le\varepsilon
\]
for rational $x,y$ in that interval.  This is uniform continuity on the
specified interval, with usable data.  A different interval may have a
different certificate.
\end{definition}

Together with point evaluation, this produces small output enclosures for
small input intervals: evaluate at a rational sample and add the continuity
error on both sides.  Conversely, shrinking interval-image bounds give
this continuity certificate.  A finite grid and coarse value enclosures
also give an absolute bound for $f$ on $[a,b]$.

This is stronger than continuity at each rational point.  We do not infer it
from compactness.  For square roots, finite algebra gives
\[
 |\sqrt{x}-\sqrt{y}|\le\sqrt{|x-y|}\qquad(x,y\ge0),
\]
so $\mu(\varepsilon)=\varepsilon^2$ works.  For a polynomial, expand
$y^k-x^k=(y-x)\sum_{j=0}^{k-1}y^{k-1-j}x^j$ and bound the finite sum on
the interval.  Separated rational functions follow by the reciprocal bound.

\begin{proposition}[Computed inputs]\label{prop:computed-inputs}
An interval-continuous rational-input computation can be evaluated on a
number computation known to stay in its interval.
\end{proposition}
\begin{proof}
Intersect each input enclosure with $[a,b]$, choose a rational sample in
that finite intersection, and request enough input precision that all its
samples have outputs within the desired continuity tolerance.  Evaluate
the sample to the remaining tolerance.  The resulting output boxes are
compatible across input refinements and shrink.  Finite-prefix intersections
make them nested.  Equivalent input computations give equivalent outputs
by the same estimate.
\end{proof}

This is the precise meaning of an implicit continuation beyond rational
inputs.  We have constructed an evaluation procedure, not postulated a
function on a completed line.

\section{Inverting a monotone computation}
\label{thm:constructive-inverse-function}
Suppose $f$ is increasing and interval-continuous on $[a,b]$, and a
positive separation bound $\sigma(d)$ is supplied:
\[
 y-x\ge d>0\quad\Longrightarrow\quad f(y)-f(x)\ge\sigma(d)>0.
\]
Let a computed target $z$ be certified between $f(a)$ and $f(b)$.
Use the two trisection points $u,v$ of the current rational bracket.
Their images are separated by $\sigma(v-u)$.  Refine the value enclosures
until either $z>f(u)$ or $z<f(v)$ is certified.  At least one strict
inequality must become visible, even when the target equals one of the
other values.  Keep $[u,b]$ in the first case and $[a,v]$ in the second.

Each step preserves a range bracket and shrinks the input interval by a
factor $2/3$.  Interval continuity proves that the resulting input
computation maps to $z$.  Separation proves uniqueness.  This is a
terminating inverse construction without a decision procedure for equality
of computed numbers.  Squaring on $[0,1]$ has
$\sigma(d)=d^2$; the angle chart in the next chapter has a linear bound.

\section{Domains and equivalent descriptions}
Functions agree only where both computations are defined.  A geometric
chart, a series disk, and a positive algebraic branch may describe the same
function on overlapping domains without being interchangeable everywhere.
Finite chains of proved comparisons retain these domain restrictions.

The rest of the book uses only the constructions just described: finite
rational operations, enclosures, terminating refinement, and equivalence.
When a new infinite process occurs, its finite future errors will be
bounded explicitly.  No theorem identifying these computations with an
ambient completed number system is required.

\chapter{The Circle, the Sphere, and the Cylinder}\label{ch:circle-sphere}

Area, length, and volume can be computed before defining an integral of a
function.  We start with finite polygons and prisms, then construct curved
figures by inner and outer comparisons whose gaps shrink.  Each particular
figure brings its own proof.  Neither trigonometric functions nor a general
integration theorem is needed in this chapter.

The sphere and cylinder results below are a coordinate reconstruction in
the spirit of Archimedes' exhaustion arguments, not a transcription of his
proofs.  His opening letters in \emph{On the Sphere and Cylinder} describe
the surface, volume, spherical-segment, and sector comparisons pursued
here.  The historical references are given at the end of the chapter.

\section{Finite rational geometry}
\label{thm:rational-circle-pythagorean}
The rational chart
\[
 P(t)=\left(\frac{1-t^2}{1+t^2},\frac{2t}{1+t^2}\right)
\]
satisfies $P_1(t)^2+P_2(t)^2=1$ by expansion.  For $0\le t\le1$ it
runs through the first quadrant, from $(1,0)$ to $(0,1)$.  Conversely a
point of the unit circle in this quadrant has parameter $t=y/(1+x)$.

Determinants give triangle and polygon areas.  Multiplying a polygon's
area by a rational height gives the volume of a right prism: triangulate
the base and use finite additivity.  Square-root computations give segment
lengths and hence the areas of triangles and trapezoids in rational
coordinate planes.  Computed lengths are enclosed by the rational
arithmetic of Chapter~\ref{ch:foundations}; they are not assumed to be
exact rational numbers.

Reflections and the coordinate matrix
\[
 \begin{pmatrix}c&-s\\s&c\end{pmatrix},\qquad c^2+s^2=1,
\]
preserve squared lengths and determinants by finite algebra.  Here $c,s$
are coordinates of a point, not previously defined trigonometric functions.
The broken-line inequality compares polygonal lengths.

\section{Archimedes and rational circles}
\label{thm:geometric-pi-validity}
Join finitely many points $P(t_i)$ to the origin, where
$0=t_0<\cdots<t_N=t\le1$.  For $u=t_i,v=t_{i+1}$, the inner triangle
and the outer quadrilateral cut out by the endpoint tangents have areas
\[
 I(u,v)=\frac{(v-u)(1+uv)}{(1+u^2)(1+v^2)},\qquad
 O(u,v)=\frac{v-u}{1+uv}.
\]
The formulas follow from determinants and the two tangent equations
$P(u)\mathbin{\cdot}z=P(v)\mathbin{\cdot}z=1$.
Inserting a point increases the total inner area and decreases the total
outer area: the changes are areas of finite triangles.  Also
\[
 \frac{v-u}{1+v^2}\le I(u,v)\le O(u,v)
       \le\frac{v-u}{1+u^2}.
\]
Write $\delta=\max_i(t_{i+1}-t_i)$.  The total outer--inner gap is at
most
\[
 \delta\sum_i\left(\frac1{1+t_i^2}-\frac1{1+t_{i+1}^2}\right)
 =\delta\left(1-\frac1{1+t^2}\right)\le\delta.
\]
Nested rational subdivisions therefore define a sector-area computation
$A(t)$.  Define
\[
 \pi=4A(1).
\]
A square and finer rational inscribed and circumscribed polygons give
$2<\pi<4$.  Each strict inequality is witnessed by a finite computation.
Scaling coordinates by $R$ scales every polygonal area by $R^2$; thus the
disk of radius $R$ has area computation $\pi R^2$.

\begin{figure}[htbp]\centering
\rationalCircleSubdivisionAnimation
\caption{Inner and outer rational polygons compute the circle's area.}
\label{fig:rational-parameter-line}
\end{figure}

For later use, the same inequalities applied to a subarc give
\[
 \frac{v-u}{1+v^2}\le A(v)-A(u)\le\frac{v-u}{1+u^2}
 \qquad(0\le u<v\le1).
\]
In particular $(v-u)/2\le A(v)-A(u)\le v-u$.  These are geometric
area estimates.  Their later interpretation as rectangle bounds is not
needed to construct them.

\section{Circumference and sector length}
Area and circumference are initially independent computations.  Approximate
the boundary by inscribed chords and circumscribed tangent segments.
For a unit-circle chord of length $c$, its midpoint is at distance
$d=\sqrt{1-c^2/4}$ from the origin.  Twice its radial triangle area is
$cd$, and
\[
 0\le c-cd\le c^3/4.
\]
If $c_*$ is the largest chord, the difference between the inner perimeter
and twice the inner area is at most $c_*^2$ times the inner perimeter
divided by $4$.  The perimeters have a fixed rational bound supplied by
an outer square.  The difference therefore shrinks.

For a polygon tangent to the unit circle, twice its area equals its
perimeter, triangle by triangle.  The inner and outer areas have already
been proved to agree.  The polygonal length computations consequently agree
as well, and give circumference $2\pi$.  Scaling gives circumference
$2\pi R$ at radius $R$.

The same comparison on a sector says that its arc length is twice its
area divided by its radius.  Additivity and invariance under rotation
follow by finite polygonal refinements and the determinant identities.
Computed coordinates are handled by rational enclosures and their finite
error bounds.  This proves the facts about arcs that the next chapter
will use to introduce angle parameters.

\section{Cylinders and cones by finite solid comparisons}
\label{sec:cone-exhaustion}
Extruding the circle's inner and outer polygons through a height $H$
gives inner and outer prisms.  Their volume gap is $H$ times the planar
area gap.  Thus a circular cylinder has volume computation
\[
 V_{\mathrm{cyl}}=\pi r^2H.
\]
Extruding boundary polygons instead gives lateral-area computation
$2\pi rH$.  Including both end disks gives total area
$2\pi rH+2\pi r^2$.

To compute a cone with apex at height zero, base radius $r$, and height
$H$, cut its axis at $z_k=kH/N$.  In slab $k$, the cone's radius lies
between $rk/N$ and $r(k+1)/N$.  Inscribed and circumscribed cylinder
stacks therefore have volumes
\[
 C_N^- =\frac{\pi r^2H}{N^3}\sum_{k=0}^{N-1}k^2,
 \qquad
 C_N^+ =\frac{\pi r^2H}{N^3}\sum_{k=1}^{N}k^2.
\]
These are finite solid containments, not a theorem about infinitesimal
slices.  Each circular cylinder can itself be replaced by polygonal prisms
with an arbitrarily small additional error.

The identity
$\sum_{k=1}^{N}k^2=N(N+1)(2N+1)/6$ follows by induction (or by
summing successive cubes).  Consequently
\[
 C_N^\pm=\pi r^2H\left(\frac13\mathbin{\pm}\frac1{2N}
                              +\frac1{6N^2}\right).
\]
Here the minus sign gives $C_N^-$ and the plus sign gives $C_N^+$.
Their gap is $\pi r^2H/N$.  Refining, for example with $N=2^n$, gives
nested solid bounds, and proves
\[
 \boxed{V_{\mathrm{cone}}=\frac13\pi r^2H.}
\]
All occurrences of $\pi$ in finite bounds are evaluated by its existing
rational interval computation.  The displayed expressions abbreviate
finite arithmetic on these enclosures.

\section{The hemisphere and the cone fill the cylinder in the bounds}
\label{sec:sphere-volume-exhaustion}
Fix a positive rational radius $R$.  Use height $z$ above the equatorial
plane.  The upper hemisphere has squared sectional radius $R^2-z^2$.
Compare it with the cone whose sectional radius is $z$, both inside the
cylinder of radius and height $R$.

For a rational partition $0=z_0<\cdots<z_N=R$, put
$\Delta z_i=z_{i+1}-z_i$.  Cylinder stacks bounding the hemisphere have
volumes
\[
 H_P^-=\pi\sum_i(R^2-z_{i+1}^2)\Delta z_i,
 \qquad
 H_P^+=\pi\sum_i(R^2-z_i^2)\Delta z_i.
\]
The corresponding cone stacks have volumes
\[
 C_P^-=\pi\sum_i z_i^2\Delta z_i,
 \qquad
 C_P^+=\pi\sum_i z_{i+1}^2\Delta z_i.
\]
At every finite partition there are exact identities
\[
 \boxed{H_P^-+C_P^+=\pi R^3,
        \qquad H_P^++C_P^-=\pi R^3.}
\]
Moreover their common gap is at most
\[
 \pi\max_i\Delta z_i\sum_i(z_{i+1}^2-z_i^2)
 =\pi R^2\max_i\Delta z_i.
\]
Thus the independently bounded hemisphere and cone volumes add to the
cylinder volume.  The cone calculation gives the hemisphere volume
$2\pi R^3/3$, and reflection gives
\begin{theorem}[Sphere volume]\label{thm:sphere-volume}
The sphere's solid-volume computation is
\[
 \boxed{V_R=\frac43\pi R^3.}
\]
It is two-thirds of the volume of the circumscribed cylinder of radius
$R$ and height $2R$.
\end{theorem}

For equal $N$-slab subdivisions on each hemisphere the explicit bounds are
\[
 \pi R^3\left(\frac43-\frac1N-\frac1{3N^2}\right)
 \le V_R\le
 \pi R^3\left(\frac43+\frac1N-\frac1{3N^2}\right).
\]
Their difference is $2\pi R^3/N$.  This is the entire exhaustion error,
apart from requested finite precision in the disk-area computations.
No general principle asserting equality of volumes from equality of
cross-sections has been invoked.

\section{A conical frustum's area is another polygonal computation}
\label{lem:frustum-area}
Rotate a segment with endpoint distances $r,r'$ from the axis and axial
height $h$.  Its length is $\ell=\sqrt{h^2+(r-r')^2}$.
We compute the lateral area of the resulting frustum from planar facets,
before measuring the sphere's surface.

Choose an inner polygon of the unit circle and the outer polygon given by
its endpoint tangents.  Let their areas be $a^-,a^+$, their perimeters
$p^-,p^+$, and let the inner chord lengths be $b_j$ with midpoint
distances $d_j=\sqrt{1-b_j^2/4}$.  Scale these two polygons by $r$ and
$r'$ in the two endpoint planes and join corresponding edges.

An inner trapezoid has area
\[
 T_j=\frac{r+r'}2 b_j\sqrt{h^2+d_j^2(r-r')^2}.
\]
Its slant height lies between $d_j\ell$ and $\ell$.  Thus its total
area $F^-$ satisfies
\[
 (r+r')\ell a^-
 =\frac{r+r'}2\ell\sum_j b_jd_j
 \le F^-\le\frac{r+r'}2\ell p^-.
\]
Every outer polygon edge has distance $1$ from the axis in its base
plane.  The transverse slant height of the associated outer trapezoid is
therefore exactly $\ell$.  Its total area is
\[
 F^+=\frac{r+r'}2\ell p^+=(r+r')\ell a^+.
\]
The planar comparisons $p^-\le2\pi$ and $a^-\le\pi\le a^+$ imply
\[
 F^-\le\pi(r+r')\ell\le F^+,
 \qquad F^+-F^-\le(r+r')\ell(a^+-a^-).
\]
Refinement constructs a unique lateral-area computation and proves
\[
 \boxed{\operatorname{Area}(\text{frustum})=\pi(r+r')\ell.}
\]
The case $r'=0$ is the cone's lateral area $\pi r\ell$; the case $r=r'$
is the cylinder's lateral area.  Degenerate zero-area cases are read
directly.  This proof used only planar trapezoids and the circle's
exhaustion, not an unproved formula for a surface of revolution.

\section{The finite chord identity behind the sphere's surface}
\label{lem:spherical-band-chord}
Use coordinates $(r,z)$ in a meridian half-plane, with $r\ge0$.
Take consecutive points
\[
 p=(r,z),\qquad q=(r',z'),\qquad
 r^2+z^2=(r')^2+(z')^2=R^2,\qquad z>z'.
\]
Exclude the single diameter joining both poles.  Set
\[
 \Delta z=z-z',\quad \ell=|p-q|,\quad
 m=(p+q)/2,\quad d=|m|=\sqrt{R^2-\ell^2/4}>0.
\]
The midpoint vector is perpendicular to the chord.  Its radial coordinate
is $(r+r')/2$, so elementary similar right triangles give
\[
 \frac{r+r'}2\ell=d\Delta z.
\]
Rotating the chord gives an inner frustum of area
\[
 L(p,q)=\pi(r+r')\ell=2\pi d\Delta z.
\]

For an outer comparison, take the parallel tangent to the circle at
$Rm/d$ and truncate it by the same two horizontal planes.  Its endpoint
radii exceed $r,r'$ by the same amount
\[
 k=(R-d)\frac{\ell}{\Delta z}.
\]
The slant length is still $\ell$, so the tangent frustum has area
\[
 U(p,q)=2\pi\left(d\Delta z+(R-d)\frac{\ell^2}{\Delta z}\right).
\]
These are actual finite frustum computations.  Since $d\le R$ and
$\ell\ge\Delta z$, their comparison is
\[
 \boxed{L(p,q)\le2\pi R\Delta z\le U(p,q).}
\]
The middle quantity is the independently computed area of the cylinder
band of radius $R$ between the same planes.

There is also a uniform error bound, including near the poles.  The
nonnegative radii give
\[
 (r-r')^2\le|r^2-(r')^2|\le2R\Delta z,
 \qquad \ell^2\le4R\Delta z.
\]
Using $R-d=\ell^2/(4(R+d))$ therefore yields
\[
 \boxed{U(p,q)-L(p,q)
       =\frac{\pi\ell^4}{2(R+d)\Delta z}\le2\pi\ell^2.}
\]
This explicit estimate, not an assertion about arbitrary triangulations
of a surface, will prove convergence.

\section{Sphere and zone surface areas by exhaustion}
\label{sec:sphere-area-exhaustion}
For a full meridian take, on its northern half,
\[
 (r(t),z(t))=
 R\left(\frac{2t}{1+t^2},\frac{1-t^2}{1+t^2}\right),
 \qquad t=j/N,\quad 0\le j\le N,
\]
and reflect in the equatorial plane for the southern half.  These are
rational points.  Adjacent northern chord lengths satisfy the exact identity
\[
 \ell(u,v)^2=
 \frac{4R^2(v-u)^2}{(1+u^2)(1+v^2)}\le\frac{4R^2}{N^2}.
\]
The reflected chords have the same bound.  Write $\ell_*$ for the largest
chord, so $\ell_*\le2R/N$.  Their total length is at most the meridian
semicircle length $\pi R$, already computed from polygons.

Let $L_P$ be the sum of inner frustum areas and $U_P$ the sum of the
separate tangent-frustum areas.  The latter need not join along a common
parallel; its validity is the finite comparison below, not a general claim
that containment of surfaces orders their areas.  Adding the chord
identities and the error estimates gives
\[
 L_P\le4\pi R^2\le U_P,
 \qquad U_P-L_P\le2\pi\ell_*\sum_i\ell_i
                  \le2\pi^2R\ell_*\le\frac{4\pi^2R^2}{N}.
\]
Thus $64R^2/N$ is a rational upper bound for the geometric gap.  All these
intervals overlap, independently of the chosen meridian subdivisions.
Compute the finitely many lengths, tangents, and areas to any further
prescribed rational error and take finite-prefix intersections.  This
constructs the sphere's surface-area computation, independently of its
volume, and proves
\begin{theorem}[Sphere area]\label{thm:sphere-area}
The sphere's surface-area computation is
\[
 \boxed{\mathcal S_R=4\pi R^2.}
\]
It equals the lateral area of the circumscribed cylinder, and equals
four times the area of a great-circle disk.
\end{theorem}

At the coarsest stage, the meridian consists of the two chords from the
north pole to the equator and then to the south pole.  Its bounds are
\[
 2\sqrt2\,\pi R^2\le\mathcal S_R
       \le(8-2\sqrt2)\pi R^2.
\]
Further rational subdivisions improve these bounds by the proved schedule.

The same proof works between any two fixed parallel planes with heights
$z_0<z_1$.  The chord heights telescope to $H=z_1-z_0$ instead of $2R$.
One may take rational heights and compute the endpoint radii by square
roots.  With height mesh $\delta$, the earlier bound gives
$\ell_*\le\sqrt{4R\delta}$, so refinement terminates even at a pole.
The construction proves
\[
 \boxed{\operatorname{Area}(\text{spherical zone of height }H)
                  =2\pi RH.}
\]
In particular a cap of height $h$ has area $2\pi Rh$.  The squared
straight-line distance from its pole to a point of its rim is $2Rh$.
The cap therefore has the area of the planar disk with that chord as
radius, another of Archimedes' comparisons.

Adding the two end disks to the circumscribed cylinder gives total area
$6\pi R^2$.  Thus both the sphere's volume and its surface area are
two-thirds of those of the cylinder \emph{including its bases}.  Lateral
area alone instead equals the sphere's area.  These are distinct statements.

\section{Caps and solid sectors}
\label{sec:spherical-caps-sectors}
The same finite sum identities also compute a spherical cap's volume.
Let $0\le h\le2R$ be rational and measure depth $t$ from the north pole.
The squared section radius is $q(t)=2Rt-t^2$.  For $t,u\in[0,2R]$,
\[
 |q(t)-q(u)|=|t-u|\,|2R-t-u|\le2R|t-u|.
\]
On $N$ equal depth slabs of thickness $\delta=h/N$, enclose squared
radii between $\max(0,q(j\delta)-2R\delta)$ and
$q(j\delta)+2R\delta$.  The associated cylinder stacks give volume
bounds whose gap is at most $4\pi R h\delta$.

Their reference sum, obtained using $q(j\delta)$, is exactly
\[
 \pi\delta\sum_{j=0}^{N-1}q(j\delta)
 =\pi\left[R h^2\left(1-\frac1N\right)
   -h^3\left(\frac13-\frac1{2N}+\frac1{6N^2}\right)\right].
\]
This uses only the finite sums of integers and squares.  Its difference
from $\pi(Rh^2-h^3/3)$ is bounded by
$\pi(Rh^2/N+h^3/(2N)+h^3/(6N^2))$, while either bounding stack differs
from the reference sum by at most $2\pi R h\delta$.  These bounds
construct the cap volume and prove
\[
 \boxed{V_{\mathrm{cap}}=\pi h^2\left(R-\frac h3\right).}
\]
For $h=0$ use the zero computation.  For $h=2R$ it agrees with the earlier,
independent full-sphere computation.

For $0\le h\le R$, the solid sector joining the cap to the center is
the cap together with a cone of height $R-h$ and squared base radius
$2Rh-h^2$.  The cone formula and finite additivity give
\[
 V_{\mathrm{sector}}
 =\pi h^2(R-h/3)+\frac\pi3(2Rh-h^2)(R-h)
 =\frac23\pi R^2h
 =\frac R3\operatorname{Area}(\text{cap}).
\]
For $R<h\le2R$, subtract the opposite smaller sector from the sphere;
the same formula results.  Taking the whole sphere gives
$V_R=R\mathcal S_R/3$, now as a consequence of two independent geometric
computations rather than a definition of either quantity.

\section{What precedes calculus}
Only finitely many polygons, prisms, cylinders, and planar frusta occur at
any stage.  The operations are determinants, square roots, rational sums,
and interval arithmetic.  Inner--outer gaps supply the needed number
computations.  The first cone and sphere sums have been evaluated by finite
integer identities, not by antiderivatives.  Surface area has its own
facet construction, not the rule of differentiating volume with respect
to radius.

The next chapter introduces trigonometry by using sector area as an angle
coordinate.  Only afterwards do we define a general rectangle integral
and identify these particular geometric exhaustions as examples of it.
That later organization explains and reuses the geometry; it is not a
prerequisite for the results proved here.

\begin{remark}[Historical source]
Archimedes, \emph{On the Sphere and Cylinder}, Books I and II, in T. L.
Heath's edition of \emph{The Works of Archimedes} (1897), especially the
opening letters stating the principal results.  The texts are available
as \href{https://doi.org/10.1017/CBO9780511695124.010}{Book I} and
\href{https://doi.org/10.1017/CBO9780511695124.011}{Book II}.
The rational-coordinate schedules and explicit error bounds in this
chapter are our reconstruction; they should not be attributed verbatim
to Archimedes.
\end{remark}

\chapter{Angles and Trigonometric Functions}\label{ch:rational-circle-trigonometry}

The preceding chapter constructed circles, arc lengths, sector areas, and
the sphere--cylinder comparisons without trigonometric functions.  We now
use sector area to assign an angle coordinate to a point of the circle.
The new objects are functions of that coordinate:
\[
 S(x)=\sin(\pi x),\qquad C(x)=\cos(\pi x).
\]
These names refer to the computations below, not to functions imported
with an existing theory of differentiation.

\section{The sector-area coordinate}
\label{sec:sector-angle-coordinate}
Recall the rational circle chart and its sector-area computation from
Chapter~\ref{ch:circle-sphere}:
\[
 P(t)=\left(\frac{1-t^2}{1+t^2},\frac{2t}{1+t^2}\right),
 \qquad A(1)=\pi/4.
\]
A full turn has normalized parameter length $2$ and a quarter turn has
length $1/2$.  For rational $0\le x\le1/2$, the point with sector area
$\pi x/2$ is to have coordinates $(C(x),S(x))$.

The geometric estimates already proved are
\[
 \frac{v-u}{1+v^2}\le A(v)-A(u)\le\frac{v-u}{1+u^2},
 \qquad \frac{v-u}{2}\le A(v)-A(u)\le v-u
\]
for rational $0\le u<v\le1$.  The chart therefore has interval
continuity and a linear inverse-separation bound.  The target
$\pi x/2=2xA(1)$ is enclosed by its endpoint values.  The trisection
construction of Chapter~\ref{ch:foundations} computes a unique parameter
$t$ such that
\[
 A(t)=\pi x/2.
\]
Evaluation of $P$ at this computed input gives $C(x),S(x)$.  At $x=0$
and $x=1/2$ use $t=0,1$ directly.  Reflections in the coordinate axes and
whole turns give the computation at every rational input.

The conventional name $A(t)=\arctan t$ records the half-angle nature of
this rational chart: the arc length from $(1,0)$ to $P(t)$ is $2A(t)$.
No integral definition of arctangent is needed.  The rational rectangle
bounds for $1/(1+t^2)$ that occur in the sector-area construction will
later give an equivalent integral computation.

\section{Dyadic rotations as another computation}
\label{thm:finite-quarter-turn}
If $(c,s)$ is a unit-circle point in the first quadrant, its half-angle
point is
\[
 (a,b)=\left(\sqrt{(1+c)/2},\sqrt{(1-c)/2}\right).
\]
Finite algebra gives $a^2+b^2=1$, $a^2-b^2=c$, and $2ab=s$; the
nonnegative sign in the last identity follows from $s\ge0$.  Thus the
coordinate rotation matrix of $(a,b)$ squares to that of $(c,s)$.
Additivity and invariance of sector area show that its angle coordinate
is half that of the parent point.

Repeated halving of the quarter turn, followed by finite rotation products,
therefore constructs all dyadic angle samples by nested square roots.
They agree with the inverse-sector construction by its separation and
uniqueness.  This agreement is a geometric theorem between two computations,
not an appeal to an already completed trigonometric function.

\section{Addition laws and continuity}
Multiplying the coordinate rotation matrices and using the additivity of
sector area gives
\begin{align*}
 S(x+y)&=S(x)C(y)+C(x)S(y),\\
 C(x+y)&=C(x)C(y)-S(x)S(y),\\
 C(x)^2+S(x)^2&=1.
\end{align*}
The endpoint values are
$S(0)=0$, $C(0)=1$, $S(1/2)=1$, and $C(1/2)=0$.
Reflections give $S(-x)=-S(x)$ and $C(-x)=C(x)$; both have period $2$.

A coordinate difference is bounded by chord length, and chord length by
arc length.  The arc-length computation from Chapter~\ref{ch:circle-sphere}
and the bound $\pi<4$ yield
\[
 |S(x)-S(y)|\le4|x-y|,\qquad
 |C(x)-C(y)|\le4|x-y|.
\]
These hold for every rational pair, splitting across quadrant boundaries
when necessary.  Thus dyadic approximation is another terminating way to
evaluate any rational angle.  They also give evaluation at computed inputs
on specified bounded intervals.

\section{Small angles before differentiation}
\label{sec:geometric-small-angle}
For $0<h\le1/4$, chord, sector, and tangent comparisons give
\[
 S(h)\le\pi h\le S(h)/C(h).
\]
For the left inequality, the vertical displacement is at most the chord,
which is at most the arc length.  For the right inequality, twice the
sector area is at most twice the area of the enclosing tangent triangle.
The geometric point lies between the positive horizontal axis and its
bisector, so $C(h)>0$ with a positive computable lower bound.

The half-angle identity and chord bound give
\[
 0\le1-C(h)=2S(h/2)^2\le8h^2.
\]
Multiplying the upper tangent comparison by $C(h)$ and using $\pi<4$
yields
\[
 0\le\pi h-S(h)\le\pi h(1-C(h))\le32h^3.
\]
Odd and even symmetry therefore give, for $|h|\le1/4$,
\[
 |S(h)-\pi h|\le32|h|^3,\qquad |1-C(h)|\le8h^2.
\]
The second estimate also holds for all rational $h$, by the global chord
bound and half-angle identity.  These finite geometric errors will prove
both that the convex derivative computations terminate and that their
values are $\pi C$ and $-\pi S$.

\section{Input to integration}
One of the first areas computed in two ways will be
\[
 \int_0^{1/2}S(x)\,dx.
\]
Equal-angle samples use the nested-radical dyadic construction.  Unequal
samples use the rational circle parameter $t$ and the sector coordinate
$2A(t)/\pi$.  Their comparison uses the explicit continuity bounds and
finite refinements.  The later derivative computation and FTC identify
the answer as $1/\pi$.

The distinction between the chapters is now explicit: the preceding one
computes particular geometric magnitudes, while this one constructs angle
functions and compares their descriptions.  General integration organizes
both sorts of computations only after they have been established.

\chapter{Integrals and the Fundamental Theorem}\label{ch:integrals}

An integral is a rational-box algorithm, not an ambient area supplied by
completeness.  Its motivating calculation is the area below the arctangent
kernel.  The same design handles every admitted special function.

\begin{verbatim}
partition -> certified value ranges on cells -> add rectangles
          -> nested boxes -> shrink-width certificate
\end{verbatim}

The point is to relate all such computations, crude or fast. Lebesgue measure
and unrestricted measurable functions are outside this foundation.

\section{Rectangles}\label{sec:integrals-ftc-contract}

For a partition \(a=x_0<\cdots<x_m=b\), cell ranges \([L_k,U_k]\) give
\[
 \sum_k (x_{k+1}-x_k)L_k\ \leq\ \int_a^b f\ \leq
 \sum_k (x_{k+1}-x_k)U_k.
\]
Left, right, and trapezoid sums are finite variants of this operation; their
comparison is rational interval arithmetic, not a limit theorem.

\begin{figure}[!htbp]\centering
\arctanRectangleEnclosureAnimation
\caption{A finite lower/upper rectangle computation.}
\label{fig:arctan-rectangle-enclosure}\end{figure}

\begin{verbatim}
sum := [0,0]
for each cell [x_k,x_(k+1)]:
    sum := sum + width(cell) * certified_range(f, cell, stage)
return sum
\end{verbatim}

\lean{ComputableAnalysis.riemannLeftInterval_add}
\lean{ComputableAnalysis.riemannRightInterval}
\lean{ComputableAnalysis.riemannTrapezoidInterval}
\lean{ComputableAnalysis.riemannTrapezoidInterval_eq_half_add_left_right}
\lean{ComputableAnalysis.riemannTrapezoidInterval_overlap_of_samples}
\leanok

\section{Validity from finite regularity}

A valid integral raw has ordered, nested boxes whose widths eventually fit
every positive rational tolerance. The certificate comes from explicit
regularity data, never from ``all continuous functions are integrable.''

The first reusable theorem accepts a rational Lipschitz bound. On dyadic cells
its width is bounded by \(2L/2^n\):

\begin{verbatim}
rational evaluator + Lipschitz L -> dyadic cell boxes
                                -> width <= 2*L/2^n -> valid raw
\end{verbatim}

\lean{ComputableAnalysis.Integral.LipschitzOnUnit}
\lean{ComputableAnalysis.LipschitzDyadic.compute_nested}
\lean{ComputableAnalysis.LipschitzDyadic.compute_width}
\lean{ComputableAnalysis.LipschitzDyadic.raw_valid}
\lean{ComputableAnalysis.LipschitzDyadic.raw_equiv_of_lipschitz_bounds}
\leanok

For monotone functions, endpoint ranges provide the boxes. A function with
finitely many certified turns is split into monotone pieces. Thus \(|x|\) is
two affine integrals, and \(x|x|/2\) gives its first finite-turn FTC identity.

\lean{ComputableAnalysis.Integral.MonotoneDarbouxSchedule}
\lean{ComputableAnalysis.Integral.monotoneDarbouxScheduleRaw_valid}
\lean{ComputableAnalysis.Integral.PiecewiseMonotoneConstructionFor}
\lean{ComputableAnalysis.Integral.piecewiseMonotoneIntegralFor_equiv_finiteRawSum}
\lean{ComputableAnalysis.Integral.absOnUnit_piecewise_primitiveFTC}
\leanok

\section{The effective FTC}

The central contract is finite telescoping:

\begin{verbatim}
derivative box on a cell -> bound its area -> sum bounds
                        -> telescope primitive increments -> F(b)-F(a)
\end{verbatim}

It is an interface a particular primitive earns. The project does not assert
a universal FTC merely because a function has a classical derivative. Affine
and polynomial witnesses are the regression base; arctangent, exponential,
and the tangent chart are non-polynomial instances.

\lean{ComputableAnalysis.DerivativeBoundFTC}
\lean{ComputableAnalysis.DerivativeBoundFTC.boundedIntegral_equiv_endpointDifference}
\lean{ComputableAnalysis.ExactFunction.affineFTCExact}
\lean{ComputableAnalysis.Integral.squareIntegralEffectiveFTC_integral_equiv_one_third}
\lean{ComputableAnalysis.ArctanEffectiveFTC.arctanEffectiveFTC}
\leanok

Convexity is useful because neighboring secants give finite derivative bounds.
Chapter~\ref{ch:effective-calculus} supplies the reusable certificates.

\section{Circle and sine}

For \(\sin(\pi x)\) on \([0,1/2]\), the ends are rational and every interior
dyadic sample has its own stage-dependent nested-radical box.

\begin{figure}[!htbp]\centering
\intervalSineIntegralStageAnimation
\caption{Dyadic cells and independently refined sine-value boxes.}
\label{fig:interval-sine-integral-stage}\end{figure}

The tangent chart gives a completed Riemann--Stieltjes computation. The
equal-dyadic nested-radical candidate is executable and checked through
stages 0--2; its remaining theorem is the all-depth geometric overlap with
the chart, not a hidden completeness step.

\lean{ComputableAnalysis.SinPiIntegral.dyadicNestedRadicalLeftSum}
\lean{ComputableAnalysis.SinPiIntegral.dyadicNestedRadicalLeftSum_width_le_of_stage}
\lean{ComputableAnalysis.SinPiIntegral.dyadicNestedRadicalLeftSum_zero_overlaps_stieltjes}
\lean{ComputableAnalysis.SinPiIntegral.dyadicNestedRadicalLeftSum_one_overlaps_stieltjes}
\lean{ComputableAnalysis.SinPiIntegral.dyadicNestedRadicalLeftSum_two_overlaps_stieltjes}
\lean{ComputableAnalysis.SinPiIntegral.DyadicTangentWitnessFamily}
\leanok

A trigonometric substitution is one finite parameter partition transported to
its rational-circle image. The Riemann--Stieltjes name emphasizes that image
cells need not be equal.

\section{Finite identities and assembly}

Substitution, integration by parts, and adjacent-interval addition begin as
finite mesh identities. They become integral identities only when both raw
computations have validity and overlap certificates.

\begin{verbatim}
two finite mesh presentations -> reindex / telescope -> overlap edge
\end{verbatim}

\lean{ComputableAnalysis.IntegralIdentities.coordinateTimesArctanForwardTwoStageMonotoneDefiniteIdentity}
\lean{ComputableAnalysis.IntegralIdentities.coordinateTimesArctanForwardTwoStageMonotoneIntegral_equiv_arctanGeom_one}
\lean{ComputableAnalysis.Logarithm.piTriangleLogReciprocalIntegral_equiv_four_arctanGeom_one}
\leanok

\section{Public interface}

To formalize a new integral, provide a rational evaluator; cell ranges or a
modulus; nesting and shrinking; and, when a value matters, an overlap with an
independent computation. For an FTC theorem add a primitive and a finite
derivative-bound telescope. That is the admission rule.

\chapter{The FTC Bridge: Effective Calculus}\label{ch:effective-calculus}

This chapter supplies the finite certificates consumed by the FTC contract; it
is not a second integral theory or a universal completeness theorem.
\begin{verbatim}
primitive + derivative boxes + mesh -> finite telescope -> endpoint overlap
\end{verbatim}
\lean{ComputableAnalysis.DerivativeBoundFTC}
\lean{ComputableAnalysis.ScheduledIntervalRegularOn}
\leanok

\section{Regression ladder}
Affine functions close in one rectangle. The square is the representative
polynomial derivative-bound schedule: \(\int_0^1x^2=1/3\). Its finite
integral and endpoint evaluators overlap; other routine degrees are obtained
by the generic monomial mechanism when needed.
\lean{ComputableAnalysis.ExactFunction.affineFTCExact}
\lean{ComputableAnalysis.Integral.squareIntegralEffectiveFTC_integral_equiv_one_third}
\leanok

\section{Product and trigonometric transport}
The product witness \(F(x)=x\arctan x\) is the first non-polynomial example.
\lean{ComputableAnalysis.IntegralIdentities.coordinateTimesArctanForwardTwoStageFTC}
\leanok

The tangent-chart square integral is complete. The remaining \(\sin(\pi x)^2\)
theorem is an all-depth geometric transport from nested-radical dyadic boxes
to tangent witnesses, including the reciprocal-\(\pi\) normalization. It is
an explicit finite-certificate target, not a missing classical existence proof.
\lean{ComputableAnalysis.SinPiIntegral.ArctanSinPiConstruction.halfIntegral_equiv_reciprocalPi_of_halfAngle_certificate_family}
\leanok

New theorems add an evaluator, finite range/derivative estimate, stage
schedule, validity proof, and one equivalence edge to the representation
chain; pairwise equivalence of every implementation is unnecessary.

\chapter{Series and Finite Remainders}\label{ch:infinite-series}

An infinite sum is useful here only when a finite prefix comes with a
bound on every finite future correction.  This produces a number
computation by stabilization.  Algebraic identities between prefixes,
combined with tail bounds, then prove identities between the computations.

\section{Geometric and alternating sums}
\label{thm:finite-geometric-series-interface}
For rational $|r|<1$,
\[
 \left|\sum_{k=N+1}^{M}r^k\right|
       \le\frac{|r|^{N+1}}{1-|r|}\qquad(M>N).
\]
The right side is a rational error tending to zero.  Thus the geometric
series computes $1/(1-r)$ by its finite identity.

If $a_n$ are nonnegative, decrease to zero with an effective bound, then
even and odd partial sums of $\sum(-1)^na_n$ interleave; pairing adjacent
terms proves that the next term bounds the remainder.  Computed rather
than rational terms are evaluated to a finite total error budget.

\begin{example}[Divergence is also a computation]
\label{thm:harmonic-dyadic-lower-bound}
Grouping the terms with $2^{j-1}<n\le2^j$ gives
\[
 \sum_{n=1}^{2^m}\frac1n\ge1+\frac m2.
\]
For every rational height, this produces an explicit finite partial sum
above it.  No infinite value is introduced.
\end{example}

The sum of reciprocal squares has the independent tail bound
\[
 \sum_{n=N+1}^{M}\frac1{n^2}
 \le\sum_{n=N+1}^{M}\left(\frac1{n-1}-\frac1n\right)\le\frac1N.
\]
This constructs the number before identifying it with $\pi^2/6$.
The identification will be a Fourier computation in
Chapter~\ref{ch:fourier}.

\section{Arctangent: a finite identity integrates to a series}
\label{thm:leibniz-arctan}
The exact polynomial identity
\[
 \frac1{1+t^2}=\sum_{k=0}^{N-1}(-1)^kt^{2k}
              +\frac{(-1)^Nt^{2N}}{1+t^2}
\]
has an integrable remainder.  On $[0,x]$ with $0\le x\le1$ its
absolute integral is at most $x^{2N+1}/(2N+1)$.  Finite polynomial
integration and the already proved angle-area comparison yield
\[
 A(x)=\sum_{k=0}^{N-1}\frac{(-1)^kx^{2k+1}}{2k+1}+R_N(x),
 \qquad |R_N(x)|\le\frac{x^{2N+1}}{2N+1}.
\]
At $x=1$ this is the Leibniz computation of $\pi/4$.  This is a proof of
agreement between a geometric computation and a series computation, not
an identification by notation.

\section{A power series is a prefix together with a majorant}
Suppose the coefficients $a_n$ are computable, and rational constants
$M,R>0$ give $|a_n|\le M R^{-n}$.  On $|x|\le r<R$, put $q=r/R$.
Then
\[
 F_N(x)=\sum_{n=0}^Na_nx^n,\qquad
 |F-F_N|\le\frac{M q^{N+1}}{1-q}.
\]
The inequality means a bound on all finite future tails, and so constructs
$F$.  Coefficient evaluations are made accurate enough that their finite
weighted error is smaller than a prescribed remaining budget.  The same
construction works in rational complex boxes on a smaller disk.

For the differentiated prefixes the tail is bounded by
\[
 \frac{M}{R}q^N\left(\frac{N+1}{1-q}+\frac{q}{(1-q)^2}\right).
\]
This follows by summing $nq^{n-1}$, or differentiating a finite geometric
identity and bounding its final terms.  Higher differentiated tails follow
in exactly the same finite way, with explicit polynomial factors in $n$.
One fixes the required order and writes its bound; no unrestricted exchange
of limits is presumed.

\begin{theorem}[Differentiation and integration of a supplied series]
\label{thm:finite-taylor-ftc-prefix}
On a smaller interval $[-r,r]$, the derivative is the computation with
coefficients $(n+1)a_{n+1}$.  Its integral is computed term by term, and
\[
 \int_a^b\sum_{n=0}^{\infty}a_nx^n\,dx
   =\sum_{n=0}^{\infty}\frac{a_n}{n+1}(b^{n+1}-a^{n+1}).
\]
\end{theorem}
\begin{proof}
Each finite polynomial has its exact algebraic derivative and primitive.
For all prefixes the second derivatives have the common bound
\[
 K=\frac{2M}{R^2(1-q)^3}.
\]
The finite polynomial remainder identity gives
$|F_N(y)-F_N(x)-F_N'(x)(y-x)|\le K|y-x|^2/2$ for a segment in
$[-r,r]$.  One can derive this identity by integrating the polynomial
second derivative twice; polynomial integration is just the binomial
identity and finite telescoping.  Pass to the computations using the
explicit tails for $F_N$ and $F_N'$.  This proves the derivative assertion.
Uniform tail bounds change an integral by at most its interval length
times the tail, proving the second assertion.
\end{proof}

For real coefficients, $F_N+Kx^2/2$ is convex, by the same finite polynomial
second-difference identity.  Passing its inequalities through the tails
proves that $F+Kx^2/2$ is convex.  Thus this derivative agrees with the
convex-secant computation of Chapter~\ref{ch:effective-calculus}; it is not
a competing notion of derivative.

\section{Products, coefficient identities, and Taylor errors}
For two supplied series, multiply finite prefixes.  Absolute majorants
bound the two omitted tails and their product.  Regrouping their finite
terms proves the Cauchy product identity.  The same argument in two
variables justifies the coefficient comparisons used for exponential
addition below.

A Taylor polynomial is another finite computation of the same function.
When a translated series or a finite remainder bound has been supplied,
its error is bounded by that data.  For example the uniform bound $K$
above gives the first-order error $K h^2/2$.  More generally a supplied
bound $B$ on the polynomial/series $(m+1)$st derivatives gives
$B|h|^{m+1}/(m+1)!$ by repeated finite polynomial integration and its
uniform tails.  We do not infer such a remainder theorem from pointwise
rational derivatives alone.

\section{Binomial square-root computations}
Put
\[
 c_0=1,\qquad c_{n+1}=\frac{2n+1}{2n+2}c_n
       =\frac{\binom{2n+2}{n+1}}{4^{n+1}}.
\]
The coefficients are rational and $0<c_n\le1$, so
$B(z)=\sum_{n\ge0}c_nz^n$ has a geometric tail on $|z|\le r<1$.
The recurrence gives the coefficient identity
\[
 2(1-z)B'(z)=B(z).
\]
By the product rules for series, every nonconstant coefficient of
$(1-z)B(z)^2$ has derivative zero, and its constant coefficient is $1$.
Consequently
\[
 (1-z)B(z)^2=1.
\]
For $0\le z<1$, positivity of the coefficients identifies this computation
with the positive inverse square root $(1-z)^{-1/2}$.  The argument is
coefficient algebra with tails; it does not use an ODE uniqueness theorem
for arbitrary rational-point solutions.

This example will reappear in the construction of an elliptic integral.
It also illustrates how a recurrence can supply a special function, while
an algebraic identity identifies what the computation evaluates.

\chapter{Exponential and Logarithm}\label{ch:exponential-logarithm}

These functions are useful tests of the distinction between defining a
computation and proving equivalence with another one.  We construct the
exponential by factorial prefixes, the logarithm by rectangles, and then
prove that they are inverse computations.

\section{Factorial prefixes}
For a rational bound $R\ge|x|$, set
\[
 E_N(x)=\sum_{n=0}^N\frac{x^n}{n!}.
\]
If $N+2\ge2R$, the absolute terms after the first omitted one have
successive ratio at most $1/2$.  Thus
\[
 |E-E_N|\le2\frac{R^{N+1}}{(N+1)!}.
\]
Finite-prefix stabilization constructs $E(x)$ on every bounded interval.
The differentiated series is identical, so the previous chapter proves
$DE=E$, with a uniform increment error bound on each such interval.

Finite binomial algebra in the double series gives
\[
 E(x+y)=E(x)E(y),\qquad E(0)=1,\qquad E(x)E(-x)=1.
\]
To justify the rearrangement, first restrict to finitely many terms;
the omitted terms are bounded by the product of the factorial majorants.
For $x\ge0$, the coefficients give $E(x)\ge1+x$.  The reciprocal
identity then proves positivity for negative $x$ as well.

The midpoint difference has the form
\[
 E(x+h)+E(x-h)-2E(x)
   =E(x)\bigl(E(h)+E(-h)-2\bigr)\ge0,
\]
since the final factor is a sum of nonnegative even powers.  Hence $E$ is
convex.  The series remainder proves that its convex derivative computation
terminates and has value $E$.

\section{Compound interest is another computation}
For $N\ge1$, the binomial coefficient of $x^k/k!$ in $(1+x/N)^N$ is
\[
 q_{N,k}=\prod_{j=0}^{k-1}(1-j/N)\quad(k\le N),
\]
and it is zero for $k>N$.  For $k\le N$,
$0\le1-q_{N,k}\le k(k-1)/(2N)$, by expanding the finite product or
inductively using $1-uv\le(1-u)+(1-v)$ for $u,v\in[0,1]$.
For $k>N$ the same upper bound for $1-q_{N,k}=1$ still holds.
Consequently, if $|x|\le R$ and $B_R$ is a rational upper bound for $E(R)$,
\[
 \left|(1+x/N)^N-E(x)\right|
 \le\frac1{2N}\sum_{k\ge2}\frac{k(k-1)R^k}{k!}
 \le\frac{R^2B_R}{2N}.
\]
Every infinite sum in this estimate has already been given a factorial
tail.  This proves the compound-interest representation by a finite error
budget.  For its usual positive-base interpretation also choose $N>R$.

\section{The logarithm by accumulation}
For rational $x>0$, define
\[
 L(x)=\int_1^x\frac{dt}{t}.
\]
On every positive rational interval $[a,b]$, the integrand is decreasing,
bounded, and has the algebraic continuity estimate
$|1/x-1/y|\le |x-y|/a^2$.  The integral is therefore constructed
independently of $E$.  The accumulation estimate proves
\[
 DL(x)=1/x.
\]
Ordered integral averages make $L$ concave.  Its one-sided secants converge
by the same continuity estimate, so this is also its computed derivative.

For $0<x<y$ in $[a,b]$,
\[
 \frac{y-x}{b}\le L(y)-L(x)\le\frac{y-x}{a}.
\]
This supplies continuity and strict separation.  Rational scaling of
partitions gives, for positive rational $x,y$,
\[
 L(xy)=L(x)+L(y).
\]
The same identity extends to computed positive inputs by the
computed-input construction and the displayed bounds.  In particular
$L(2^n)=nL(2)$, with $L(2)\ge1/2$, gives explicit arbitrarily large
values, and reciprocation gives arbitrarily negative ones.

\section{Proving the inverse relation}
On a fixed bounded input interval choose rational bounds $B_E,K_E$ such
that $|E|\le B_E$ and
\[
 |E(x+h)-E(x)-E(x)h|\le K_Eh^2/2.
\]
The factorial tails and the series remainder supply them.  Also enclose
the image in $[c,d]$ with rational $c>0$, using $E(-x)=1/E(x)$.
Accumulation and the Lipschitz bound $c^{-2}$ for $1/t$ give
\[
 |L(z+k)-L(z)-k/z|\le k^2/(2c^2)
\]
for a segment in that image interval.  The derivative bound gives
$|E(x+h)-E(x)|\le B_E|h|$ (enlarge $B_E$ on the interval if necessary).
Substitution into these two estimates proves the explicit composite error
\[
 |L(E(x+h))-L(E(x))-h|
 \le\left(\frac{K_E}{2c}+\frac{B_E^2}{2c^2}\right)h^2.
\]

Adding these estimates on equal rational partitions shows that
$L(E(x))-x$ has endpoint difference at most
$\varepsilon$ times the interval length for every $\varepsilon>0$.
At zero it is zero.  Hence
\[
 L(E(x))=x.
\]
This is a finite telescoping proof, not the invalid implication
``pointwise rational derivative zero implies constant.''

The identities extend to computed inputs on bounded intervals by
continuity.  In particular, for a positive computed $y$, apply the identity
to $x=L(y)$.  Both $E(L(y))$ and $y$ have logarithm $L(y)$, and the
separation bound for $L$ proves their equivalence:
\[
 E(L(y))=y.
\]
Thus $E=\exp$ and $L=\log$ name inverse computations.

\section{Representative consequences}
Define $x^\alpha=E(\alpha L(x))$ on a positively separated interval,
for a supplied computed constant $\alpha$.  The same finite increment
composition gives
\[
 D(x^\alpha)=\alpha x^{\alpha-1}.
\]
This statement includes the domain and bounded-interval estimates; it is
not an assertion that arbitrary rational-point differentiable functions
can be freely composed and integrated.

A particularly useful primitive is $xL(x)-x$.  The product calculation in
Chapter~\ref{ch:effective-calculus} proves
\[
 \int_a^b L(x)\,dx=[xL(x)-x]_a^b\qquad(0<a<b).
\]
Another is the exponential itself:
\[
 \int_a^b E(x)\,dx=E(b)-E(a).
\]
The left sides are rectangle computations and the right sides are finite
combinations of independently constructed values.

Finally define, by the factorial series, $E(z)$ for rational complex
inputs.  Coefficient separation gives
\[
 E(i\pi x)=C(x)+iS(x).
\]
For completeness, the even and odd factorial prefixes have derivatives
$C'=-\pi S$, $S'=\pi C$ and the correct initial values.  Equality with
the geometric pair follows from the following finite uniqueness argument:
the difference satisfies its integral equation, and iterating that
equation $n$ times bounds it by $B(\pi h)^n/n!$ on an interval of length
$h$.  The integral equations follow by telescoping the already proved
uniform increment bounds for both constructions.  This error tends to
zero.  Chapter~\ref{ch:differential-equations} gives the matrix form of
this same argument.

\chapter{Algebraic Branches and Rational Curves}\label{ch:algebra-fta}

An algebraic formula is useful only after a branch has been selected and
its values enclosed.  We develop representative constructions: positive
square roots, a certified simple complex root, and a rational curve.  Each
has its own finite domain and separation data.

\section{Positive roots and their derivatives}
On $[a,b]$ with rational $a>0$, the square-root bisection computation has
\[
 \frac{\sqrt y-\sqrt x}{y-x}
      =\frac1{\sqrt y+\sqrt x}\qquad(x\ne y).
\]
As both inputs increase these slopes decrease, so the function is concave.
On the interval the denominator is at least $2\sqrt a$, with an explicit
rational positive lower bound obtainable by refinement.  The difference
between the left and right secants at $x$ tends to zero by the root
continuity estimate.  Consequently
\[
 D\sqrt x=\frac1{2\sqrt x}.
\]
The derivative is not bounded near zero; no theorem requiring a bounded
endpoint derivative is applied there.

For a primitive of the square root, set $F(x)=\frac23 x\sqrt x$ for
$x\ge0$, and extend by zero for $x\le0$.  If $0\le x<y$, write
$u=\sqrt x$, $v=\sqrt y$.  Then
\[
 \frac{F(y)-F(x)}{y-x}
       =\frac{2(v^2+uv+u^2)}{3(u+v)}
\]
when $u+v>0$.  It lies between $u$ and $v$: after clearing the positive
denominator the two differences are
$(v-u)(2v+u)$ and $(v-u)(v+2u)$.  Adjacent secants are therefore ordered;
$F$ is convex and $DF(x)=\sqrt x$ on $x\ge0$, including the shrinking
one-sided bound at zero.  The convex FTC gives
\[
 \int_0^b\sqrt x\,dx=\frac23 b\sqrt b.
\]
The reciprocal-root integral can also be computed as an improper integral:
removing $[0,\varepsilon]$ has error $2\sqrt\varepsilon$, by the
primitive $2\sqrt x$ on positive interior intervals and their FTC.  The tail bound, not a bare
improper-integral symbol, defines the computation.

\section{Separated rational functions have finite remainder bounds}
Let $P,Q$ be rational polynomials on a fixed interval with the supplied
separation $|Q|\ge c>0$.  For numbers $v,w$ separated from zero,
\[
 \frac1w-\frac1v+\frac{w-v}{v^2}
       =\frac{(w-v)^2}{v^2w}.
\]
Finite polynomial expansions and this identity give
\[
 \frac{P(x+h)}{Q(x+h)}-\frac{P(x)}{Q(x)}
 =\frac{P'(x)Q(x)-P(x)Q'(x)}{Q(x)^2}h+R(x,h),
 \qquad |R(x,h)|\le K h^2,
\]
with a rational $K$ computed from polynomial coefficient bounds, $c$, and
the interval length.  Indeed the reciprocal remainder is bounded by
$c^{-3}|Q(x+h)-Q(x)|^2$, and the polynomial linear remainders are bounded
by their finite binomial expansions.  Multiplication adds only bounded
quadratic terms.

Applying the estimate at $h$ and $-h$ cancels the linear terms.  Thus
$P/Q+Kx^2$ is midpoint convex, hence convex, and its centered secant gaps
are at most $4K|h|$.  The quotient formula is consequently the actual
computed derivative, with a convex-difference FTC certificate.  This
supplies, in particular, the rational substitutions in this chapter and
the final elliptic-integral example.  No assertion about an arbitrary
pointwise rational second derivative is used.

\section{A complex root certified by contraction}
Let $P$ be a polynomial with rational complex coefficients, and let $c$
be a rational complex point with $P'(c)\ne0$.  Define
\[
 T(z)=z-P(z)/P'(c).
\]
Use the coordinate maximum norm on complex numbers.  Multiplication by
$w$ has norm at most $|\Re w|+|\Im w|$.  Suppose finite polynomial
bounds on the square $\|z-c\|_\infty\le r$ certify
\[
 \|T(c)-c\|_\infty\le(1-q)r,\qquad
 \|T(z)-T(w)\|_\infty\le q\|z-w\|_\infty,
 \quad 0\le q<1.
\]
The second inequality can be checked by factoring $P(z)-P(w)$ and
bounding the resulting polynomial on two copies of the square.
Then $T$ preserves the square.  Starting at $z_0=c$, compute the rational
iterates $z_{n+1}=T(z_n)$.  For $m>n$,
\[
 \|z_m-z_n\|_\infty
       \le\frac{q^n}{1-q}\|z_1-z_0\|_\infty.
\]
Rational boxes with this error, stabilized coordinatewise, define a root
computation.  Polynomial continuity and the iteration identity prove
$P(z)=0$.  Two root computations in the square have distance at most
$q^n$ times their initial distance for every $n$, proving uniqueness
without a fixed-point existence theorem.

\begin{example}[A nonrational complex root]
For $P(z)=z^2+2$, take $c=3i/2$ and $r=1/6$.  The residual step has norm
$1/12$.  Factoring $T(z)-T(w)$ gives contraction constant $q=2/9$ on
the square, since
\[
 T(z)-T(w)=(z-w)\left(1-\frac{z+w}{3i}\right).
\]
Also $1/12\le(1-2/9)/6$.  Thus the iteration constructs its unique root
there.  Comparison with the square-root construction identifies it as
$i\sqrt2$; reflection gives the other root.  The exact polynomial identity
$(z-i\sqrt2)(z+i\sqrt2)=z^2+2$ accounts for both roots.
\end{example}

This is a certificate for supplied root locations, not an invocation of a
general complex root-existence theorem.  Polynomial identities, separation,
and refinement are the concrete mathematical ingredients.

\section{Rational parametrization changes the computation}
The singular cubic
\[
 y^2=x^2(1-x)
\]
has the parametrization
\[
 x=1-t^2,\qquad y=t(1-t^2).
\]
Substitution verifies the equation by finite algebra.  Away from $x=0$,
the inverse parameter is $t=y/x$.  Both coordinates are rational at
rational $t$.  A rational differential form in $x,y$ therefore becomes a
rational function of $t$ wherever the denominators are separated from
zero.

For example, on $0<t<1$,
\[
 \frac{dx}{y}=-\frac{2\,dt}{1-t^2}.
\]
Between rational parameters $0<u<v<1$, its integral is computed by
\[
 -\log\frac{1+v}{1-v}+\log\frac{1+u}{1-u}.
\]
Indeed $-2/(1-t^2)=-1/(1-t)-1/(1+t)$, and the logarithm FTC proves the
formula.  The line-integral computation is defined along the parametrized
path, so its orientation and domain are explicit.

A parametrization is more than a closed formula: it can replace an
algebraic-value algorithm by exact rational evaluation and turn a difficult
integral into a previously understood one.  The final chapter will use a
rational circle chart to regularize an integral whose next useful description is a differential equation
rather than an elementary primitive.

\chapter{Local Calculations and Nonlinear Constructions}
\label{ch:local-models}

A named function should not have to repeat the foundations every time it
is differentiated, composed, or evaluated in a new region.  What can be
reused is not an unrestricted function space but a finite description
with error bounds.  This chapter develops that description for local
series, and a second construction for nonlinear inverses and evolutions.
The earlier convex derivative remains the first route to the FTC; the
present local estimates provide further concrete FTC certificates.

\section{A series chart and its finite jets}
A chart about a rational complex point $c$ supplies coefficient
computations $a_n$ and positive rational constants $M,R$ with
\[
 F(c+w)=\sum_{n\ge0}a_nw^n,\qquad |a_n|\le M R^{-n}.
\]
The displayed equality names the prefix-and-tail computation, not an
independently assumed function.  The disk means inputs with a certified
bound $|w|\le r<R$.  Complex absolute values are computed by square roots;
all required upper bounds can be rational.  Put $q=r/R$.

A finite jet of order $k$ is the list of the first $k$ derivatives at a
specified center.  Coefficients give that list by $F^{(j)}(c)=j!a_j$.
For evaluation away from the center, the derivative of order $k$ has the
explicit tail, after original coefficient $N\ge k-1$,
\[
 T_{N,k}=\frac{M}{R^k}\sum_{n>N}(n)_k q^{n-k},
 \qquad (n)_k=n(n-1)\cdots(n-k+1).
\]
Here and below $(n)_k$ in a derivative-tail formula is a falling factorial;
rising factorials will always be identified when introduced.  This tail
is computable without an infinite summation procedure.  Once
\[
 q\frac{n+1}{n+1-k}\le\theta<1,
\]
the successive terms decrease geometrically; their remaining sum is at
most the first one divided by $1-\theta$.  The ratio decreases with $n$,
so one finite threshold proves the estimate for every later term.  The
case $q=0$ is direct coefficient evaluation.  Coefficient-box errors are
allocated separately to the finitely many retained terms.

The full derivative bound is particularly simple:
\[
 |F^{(k)}(c+w)|\le\frac{M k!}{R^k(1-q)^{k+1}}.
\]
It follows by repeatedly differentiating the finite geometric identity
and bounding its terminal terms.  Thus the chart supplies derivatives of
any requested finite order, rather than asking a value algorithm to reveal
its derivatives.

\section{Moving the center}
\begin{theorem}[Recentring with a new majorant]
\label{thm:recenter-model}
Let $|d|\le d_0<R$.  About $c+d$, the same computation has coefficients
\[
 b_k=\sum_{n=k}^{\infty}\binom nk a_nd^{n-k},
 \qquad |b_k|\le\frac{MR}{(R-d_0)^{k+1}}.
\]
They form a chart with radius $R-d_0$ and constant $MR/(R-d_0)$.
\end{theorem}
\begin{proof}
For a finite polynomial the change of center is the binomial identity.
The majorant for $b_k$ is
\[
 \frac{M}{R^k}\sum_{j\ge0}\binom{k+j}{k}(d_0/R)^j
 =\frac{M}{R^k}(1-d_0/R)^{-k-1}.
\]
The equality follows from the differentiated geometric identity.  To
compute $b_k$, bound the tail by the eventual ratio
$(d_0/R)(k+j+1)/(j+1)<1$.  On $|v|+d_0<R$, the omitted double sum is
absolutely bounded by a geometric series.  Finite binomial rearrangements
therefore prove $F(c+d+v)=\sum b_kv^k$ with explicit errors.
\end{proof}

Recentring alone does not create continuation beyond a certified original
domain: its new disk is contained in that domain.  To move farther one
needs new bounds supplied by an equation, an integral, or another proved
representation.  This distinction prevents an apparent continuation
algorithm from repeatedly spending a radius it never replenishes.

\section{Arithmetic, reciprocals, and composition}
Finite polynomial arithmetic and the same majorants give sum and product
computations.  For instance, two bounds $M R^{-n}$ and $N R^{-n}$ give
product coefficients at most $MN(n+1)R^{-n}$.  At any smaller radius
$S<R$, these are at most
\[
 \frac{MN}{(1-S/R)^2}S^{-n}.
\]
Restriction costs a bound but gives another chart of the same kind.

For a reciprocal write $F=a_0+H$, supply $|a_0|\ge m>0$, and choose $r$
so that the coefficient majorant for $H$ on that disk is at most
$\theta m$, with $\theta<1$.  Then
\[
 \frac1F=\frac1{a_0}\sum_{j\ge0}(-H/a_0)^j,
 \qquad
 \left|\text{tail after }j=N\right|
 \le\frac{\theta^{N+1}}{m(1-\theta)}.
\]
The coefficients can instead be generated by the finite recurrence
\[
 b_0=a_0^{-1},\qquad
 b_n=-a_0^{-1}\sum_{j=1}^na_jb_{n-j}.
\]
The geometric error proves agreement.  There is no decision of whether an
arbitrary computed $a_0$ is zero; the separation is an explicit input.

For composition, recenter the outer chart at $F(c)$ and write
$F(c+w)=F(c)+H(w)$, where $H$ has zero constant coefficient.  Choose an
inner radius on which its absolute coefficient majorant is at most
$s<R_G$, the outer radius.  The first $N$ coefficients of
$G(F(c)+H)$ require only finitely many coefficients of $H$, since
$H^j$ has order at least $j$.  The outer tail is bounded by
\[
 M_G(s/R_G)^{N+1}/(1-s/R_G).
\]
The absolute sum of all composed coefficients weighted by the chosen inner
radius is at most $M_G/(1-s/R_G)$.  Hence each coefficient has that
geometric majorant, and degree truncation on a strictly smaller disk has
its own geometric tail, separate from the outer-series truncation.
Finite errors in evaluating $H$ are multiplied by the outer first-derivative
bound on a slightly larger certified disk.  This proves both the
composition computation and its chain rule.  Computed centers cause no
new problem: all their uses are interval operations with the displayed
strict margins.

These constructions also resolve removable expressions.  For example,
instead of dividing $S(x)$ by $x$ at zero, shift the coefficients of its
odd series.  This defines $S(x)/x$ there with value $\pi$.  A cancellation
proved in the coefficient algebra is different from a numerical division
by a small, uncertain denominator.

\section{Which approximations support which operations?}
Suppose finite computations $F_N$ approximate a function uniformly to
error $e_N$ on $[a,b]$.  Then their integral computations differ by at most
$(b-a)e_N$.  This is a finite rectangle estimate and proves integration
of the limit computation whenever the $e_N$ shrink effectively.

To differentiate, more is needed.  If $F_N'$ converge uniformly to a
computed $G$, then their finite FTC identities give
\[
 F(y)-F(x)=\int_x^yG(t)\,dt.
\]
A supplied modulus for $G$ makes this a derivative and FTC certificate
for $F$.  Uniform convergence of values alone gives no such conclusion.
The derivative tails of a series chart supply exactly the stronger data.
This is the form of a convergence theorem that later constructions use.


\section{A smooth computation which is not a single analytic chart}
\label{sec:flat-cutoff}
Define at rational inputs
\[
 B(x)=\begin{cases}E(-1/x^2),&x\ne0,\\0,&x=0.\end{cases}
\]
Rational equality decides the displayed branch.  Away from zero its
successive derivatives have the explicit form
\[
 D^nB(x)=P_n(1/x)E(-1/x^2),\qquad
 P_0(u)=1,\quad P_{n+1}(u)=2u^3P_n(u)-u^2P_n'(u).
\]
This is a finite rational polynomial recurrence.  To control the missing
point, let $d$ be the degree of $P_n$ and choose an integer $m$ with
$2m\ge d+2$.  For $0<|x|\le\delta\le1$, positivity of the exponential
series gives $E(1/x^2)\ge |x|^{-2m}/m!$.  Writing
$P_n(u)=\sum_j p_ju^j$, it follows that
\[
 |P_n(1/x)E(-1/x^2)|
 \le m!\sum_j|p_j|\,|x|^{2m-j}
 \le m!\left(\sum_j|p_j|\right)\delta^2.
\]
In particular each derivative extends continuously by zero, and its
quotient by $x$ tends to zero with a linear error bound.  Induction proves
that these extensions are the actual successive derivatives, all zero
at the origin.  This is not an inference from pointwise differentiability
on a punctured rational interval; the estimate explicitly controls the join.

Bounds on a compact interval are obtained by splitting at rational
$\pm\delta$: the estimate handles the center, while ordinary series and
reciprocal bounds handle the two separated pieces.  The same procedure
supplies moduli for every specified derivative order.  The FTC across the
join follows from the certified identities on the separated pieces and
the vanishing bounded end errors.

For rational $a<b$, set $\psi(x)=B(x-a)B(b-x)$ on $a<x<b$ and zero
outside.  All its derivatives match zero at both endpoints by the finite
product rule and the preceding estimates.  This constructs a compactly
supported smooth test with any requested derivative and error bound.
The positive value at an interior rational point shows that it is not
identically zero, although every derivative at an endpoint vanishes.
No analytic chart spanning that endpoint can represent it.  Explicit
nonanalytic functions therefore belong to the same calculus, by their
own joining estimates rather than by an assertion of analyticity.

\section{A nonlinear inverse: Lambert's function}
\label{sec:lambert-construction}
For $|z|\le1/4$, start with $w_0=0$ and compute
\[
 w_{n+1}=zE(-w_n).
\]
The exponential series gives $E(1/2)<2$ by a finite prefix and tail.
Consequently this map sends $|w|\le1/2$ into itself and has Lipschitz
constant at most $1/2$ there.  The estimate follows either from its
uniform increment identity along the segment or directly from the
exponential difference formula.  Finite iteration proves
\[
 |w_m-w_n|\le 2^{1-n}|z|\qquad(m\ge n).
\]
Successively accurate interval evaluations, with a geometric error budget
for their finite compositions, give a valid computation $W(z)$.  It
satisfies $W E(W)=z$ by the iteration identity and continuity.  Two fixed
points in this disk are equal by repeated contraction.  This constructs
one specified branch; it does not silently choose a branch elsewhere.

Parameter comparison in the fixed-point equations gives
$|W(z)-W(v)|\le4|z-v|$.  The finite exponential remainder for
$P(w)=wE(w)$ has linear coefficient $E(w)(1+w)$, whose modulus on this
disk is at least $1/4$: $|1+w|\ge1/2$ and $|E(w)|\ge1/2$.
Substitute $w=W(z)$ and $w+\Delta w=W(z+h)$ into this remainder and use
$|\Delta w|\le4|h|$.  Division by the separated linear coefficient gives
\[
 DW(z)=\frac{E(-W(z))}{1+W(z)},
\]
with an explicit quadratic increment remainder.  All constants can be
read from the exponential bounds on $|w|\le1/2$.  On real subintervals,
that remainder also supplies the convex-difference FTC certificate by
adding a sufficiently large rational multiple of $x^2$.

The example is named in the standard notation of
\href{https://dlmf.nist.gov/4.13}{DLMF, Section 4.13}; the finite iteration
and its error estimate are the definition used here.

\section{Nonlinear initial-value problems on a certified slab}
\label{sec:nonlinear-slab}
Let a specific computation $Q(t,y)$ have interval-continuity data on
$|t-t_0|\le h$, $\|y-y_0\|\le r$, with rational bounds
\[
 \|Q\|\le B,\qquad
 \|Q(t,y)-Q(t,v)\|\le L\|y-v\|,\qquad
 hB\le r,\quad q=hL<1.
\]
The vector norm is the maximum norm.  These are bounds to be proved for
the displayed formula $Q$, not consequences of a general ODE existence
theorem.  Define finite Picard computations
\[
 y_0(t)=y_0,\qquad
 y_{n+1}(t)=y_0+\int_{t_0}^t Q(u,y_n(u))\,du.
\]
Every iterate stays in the box, and finite induction gives
\[
 \|y_{n+1}-y_n\|\le hBq^n,\qquad
 \|y_m-y_n\|\le\frac{hBq^n}{1-q}\quad(m\ge n).
\]
Thus the iterates construct a uniformly approximated solution.  Passing
through the bounded interval integrals gives its integral equation;
accumulation then proves $Dy=Q(t,y)$.  Two supplied integral solutions
in the slab agree by the same contraction estimate.

For example, $y'=1+y^2$, $y(0)=0$ has $r=1$, $h=1/4$, $B=2$,
$L=2$, and $q=1/2$.  Its finite iterates are rational polynomials.
On this interval the already constructed computation
$S(t/\pi)/C(t/\pi)$ has the same equation and initial value; the cosine
is positively separated and its increment identity proves the integral
equation.  The two computations therefore agree.  Continue only through
further certified slabs; no continuation through the tangent's pole is
asserted.  Nonlinear special-function equations can be admitted in this
same concrete manner without forcing them into the linear theory.

\chapter{Complex Paths and Analytic Continuation}
\label{ch:complex-paths}

Complex analysis can begin with polynomial identities, finite paths, and
the series charts just constructed.  The contour, its subdivision, and
the local bounds are supplied data.  We prove the integral identities
actually used, rather than first postulating an arbitrary holomorphic
function on an arbitrary open set.

\section{Integrating on a segment}
For rational complex endpoints $p,q$, define
\[
 \int_{[p,q]}F(z)\,dz
 =(q-p)\int_0^1 F(p+t(q-p))\,dt.
\]
The right side consists of two real integral computations.  A polygonal
path is a finite list of such segments, and its integral is their sum.
Its length has a rational upper bound obtained from coordinate differences.
If two integrands differ by at most $\eta$ on the path of length at most
$L$, their integrals differ by at most $L\eta$.

In a single series disk, the primitive
\[
 P(c+w)=\sum_{n\ge0}\frac{a_n}{n+1}w^{n+1}
\]
has a geometric tail on every smaller disk.  Its finite polynomial
identities and their tails prove
\[
 \int_{[p,q]}F=P(q)-P(p)
\]
whenever the segment stays in such a disk.  This also holds for a
specified continuously differentiable parametrized arc in the disk: the
finite polynomial chain rule gives the result, and its error is bounded
by a supplied arc-length bound.  Circle arcs from the earlier geometric
construction are one such family.

\section{A finite form of the contour theorem}
\label{thm:finite-contour}
Suppose a rational triangulated polygon is supplied, and each closed
triangle lies in a smaller certified series disk of $F$.  The descriptions
are proved to agree on their overlaps.  The integral around each triangle
is zero, by its local primitive.  Add these finitely many identities:
interior edges cancel with their opposite orientations.  Therefore
\[
 \int_{\partial\Omega}F(z)\,dz=0.
\]
A polygon with holes is treated by the same cancellation, giving the
outer boundary integral as the sum of the positively oriented inner
boundary integrals.  A finite chain of these subdivisions proves equality
of two path integrals.  No existence of a finite subcover is assumed:
the finite subdivision and the certified disk margins are exhibited.

This is also a computational form of contour deformation.  It is
important to retain the overlap comparisons.  Merely evaluating the
same printed formula on two sides of a branch cut is not a proof that
one has continued the same function.

\section{Coefficient extraction on a circle}
For $0<r<R$, parametrize the circle by
\[
 z(t)=c+rE(2\pi it),\qquad 0\le t\le1.
\]
This is computed at rational $t$ by trigonometry, not by assuming circle
points are rational.  Its derivative and length bound have already been
proved.  Finite trigonometric orthogonality, or the trigonometric FTC,
gives for every integer $j$
\[
 \int_0^1 E(2\pi ijt)\,dt=\begin{cases}1,&j=0,\\0,&j\ne0.\end{cases}
\]
Integrate a finite polynomial and use its uniform series tail.  The result is
\[
 \boxed{a_k=\frac1{2\pi i}\int_{|z-c|=r}
                  \frac{F(z)}{(z-c)^{k+1}}\,dz.}
\]
If the value boxes on this circle certify $|F|\le B$, this formula gives
$|a_k|\le B r^{-k}$.  Thus an independent boundary estimate can produce
better coefficient bounds, rather than just use the original majorant.

For a computed $w$ with $|w-c|<r$, expand
$(z-w)^{-1}=\sum_{j\ge0}(w-c)^j/(z-c)^{j+1}$ on the circle.  Its tail
is bounded by a geometric series with ratio $|w-c|/r$.  Finite integration
and the coefficient formula give the local Cauchy identity
\[
 \boxed{F(w)=\frac1{2\pi i}\int_{|z-c|=r}
                         \frac{F(z)}{z-w}\,dz.}
\]
Both sides are now specified computations with independent truncation and
quadrature errors.  We have proved the identity for the supplied chart,
not invoked it to manufacture a chart from unverified pointwise data.

\section{Arctangent normalizes the Cauchy pole}
\label{sec:cauchy-arctan}
There is a direct rectangle route to the same normalization. Split the
counterclockwise boundary of $[-1,1]^2$ into eight affine half-edges.
On the right side, $z=1+iu$ gives
\[
 \frac{dz}{z}=\frac{u+i}{1+u^2}\,du.
\]
Pair opposite parameter signs and rotate through the four sides. The real
parts cancel exactly, leaving $8i/(1+u^2)$ on $[0,1]$. Our rational boxes
enclose every tagged sum on the common partition and have imaginary width
at most $16/2^n$. Thus the checked square computation satisfies
\[
 \oint_{\partial[-1,1]^2}\frac{dz}{z}=8iA(1)=2\pi i,
\]
with the original geometric $\pi$. This is
\texttt{PDE.CauchyContour.raw\_equiv\_twoPiI}; translation and rational
nonzero dilation preserve its pole pullback.

PNT+ formalizes the corresponding rectangle argument in Mathlib:
arctangent and logarithm side primitives normalize the pole, while
Mathlib's rectangle Cauchy theorem removes the regular part. Its checked
finite residue theorem concerns simple poles with no boundary poles.
Applying that theorem to
\[
 \frac{F(z)}{z-p}=
 \frac{F(z)-F(p)}{z-p}+\frac{F(p)}{z-p}
\]
gives the Cauchy value formula once the divided difference is extended
holomorphically across $p$. PNT+'s explicitly named derivative formula
uses a separate Mathlib circle-integral theorem. These are external
formal references, not imports into our foundation.

The native square normalization is checked. General native contour
transport, regular-part removal and the Cauchy integral formula still
require their explicit subdivision, local chart and convergence bridges.
The adjacent general contour descriptions give the mathematical program,
not additional checked declarations supplied by this normalization.
See the source ledger \texttt{docs/PNT\_CAUCHY\_ARCTAN.md}.

\section{Known poles and finite residues}
Suppose near a specified point $s$ the computation has a proved description
\[
 F(z)=\sum_{j=1}^m b_{-j}(z-s)^{-j}+H(z),
\]
where $H$ has a series chart.  On a sufficiently small certified circle,
finite coefficient extraction gives
\[
 \int F(z)\,dz=2\pi i b_{-1}.
\]
The coefficient $b_{-1}$ is called the residue.  The finite contour
cancellation above transports this result to larger supplied polygonal
contours when the region between them is covered by certified charts.
For circular boundary components, approximate the arcs by the geometric
polygonal construction inside the annular chart cover.  Uniform continuity
and path-length bounds control the integral errors; radial joining edges
cancel.  This supplies the same transport for those boundaries.

As a direct example, on $|z|=2$,
\[
 \frac1{z(z-1)}=\sum_{n\ge0}z^{-n-2},
\]
so its integral is zero.  Around $0$ and $1$ separately the residues are
$-1$ and $1$.  The two ways of computing the contour integral agree.
The poles were supplied and their expansions proved; no unsupported
search for every pole or zero has taken place.

\section{Logarithm along a specified path}
Start near $1$ with
\[
 L(1+w)=\sum_{n\ge1}\frac{(-1)^{n+1}}n w^n,\qquad |w|<1.
\]
Its derivative is the geometric series for $1/(1+w)$.  The series product
and chain rules show $E(L(1+w))=1+w$: its quotient by $1+w$ has all
nonconstant derivative coefficients zero and constant value $1$.
On the positive real interval it agrees with the previously constructed
logarithm by its primitive identity and the value at $1$.

Suppose a value $L(c)$ has been continued to a nonzero computed point $c$.
Choose a certified $d<|c|$ and set
\[
 L(c+w)=L(c)+\sum_{n\ge1}\frac{(-1)^{n+1}}n(w/c)^n,
 \qquad |w|\le d.
\]
This is a fresh chart, whose radius comes from separation from zero.
The value at the next center is computed from the old chart, and the two
charts have derivative $1/z$ on their overlap.  Their local primitive
identities and matching value prove agreement on the overlap, which is
convex.  A finite chain of such steps computes continuation along a
supplied path avoiding zero.  A rational positive distance from zero on
each path piece gives a finite step schedule.

The end value can depend on the path.  On the unit circle,
\[
 \int\frac{dz}{z}=\int_0^1 2\pi i\,dt=2\pi i.
\]
Thus continuation from $1$ once counterclockwise changes the logarithm
value by $2\pi i$.  Starting with value $0$, the computation $E(L/2)$
returns to $-1$, not $1$.  This constructs the sign change of the square
root after a circuit.  A path and a starting branch are part of the input;
there is no single globally chosen square-root computation on the
punctured plane.

\section{When does a chain of charts count as continuation?}
A continuation certificate consists of finitely many local charts, steps
lying strictly within their certified disks, and proved comparisons on
each successive overlap.  Comparisons can be coefficient identities,
local primitive identities, or the uniqueness estimates for an explicitly
constructed differential or integral equation.  Equality of all Taylor
coefficients at a common center implies equality on their common smaller
disk by the finite tail estimates.  Matching only a finite jet is not
enough unless an equation proves that it determines the later coefficients.

A path change is justified by a finite subdivision whose cells stay in
such common charts.  It does not follow from drawing apparently similar
paths.  This requirement exposes the separation and branch information
that an appeal to unrestricted analytic continuation would hide.

The next chapters use two different ways to replenish local bounds:
parameter integrals supply new majorants in right half-planes, while
linear equations supply new evolution bounds away from their coefficient
singularities.  Both can be connected to these local contour computations.

\chapter{Improper Integrals, Parameters, and the Gamma Function}
\label{ch:improper-parameters}

A specific improper integral needs bounds at its missing endpoints, not a
general theory of measurable functions.  A parameter derivative needs
corresponding bounds for its finite Taylor errors.  Gamma, beta, and the
error function give concrete tests of both requirements.

\section{Two tails and one bounded computation}
Suppose $f(t,z)$ has an integral computation on each truncated interval
$[\varepsilon,T]$, and uniform bounds on every finite end piece are supplied:
\[
 \left|\int_u^\varepsilon f(t,z)\,dt\right|\le e_0(\varepsilon),
 \quad
 \left|\int_T^V f(t,z)\,dt\right|\le e_\infty(T)
 \quad(0<u<\varepsilon<T<V).
\]
Both errors must shrink by proved schedules.  Evaluate the middle integral
to error $\eta$ and enlarge its rational box by
$e_0(\varepsilon)+e_\infty(T)$.  Comparisons on a larger common truncated
interval prove compatibility of any two such boxes.  Finite intersections
therefore construct $\int_0^\infty f(t,z)\,dt$.

These tail bounds can be absolute-integral bounds, but need not be: later
an oscillatory cancellation will give a bound without absolute convergence.
No interchange of operations is inferred merely from the notation
$\int_0^\infty$.

\section{Parameter series with an integrable majorant}
\label{thm:parameter-majorant}
Suppose on every truncated interval a supplied identity is
\[
 f(t,z_0+w)=\sum_{n\ge0}a_n(t)w^n,
 \qquad |a_n(t)|\le M(t)R^{-n},
\]
where $M$ is a nonnegative concrete integrand with proved endpoint tails
and an integral bounded by a rational $B$.  Each $a_n$ is integrable on
truncated intervals with the supplied continuity data.  The same bounds
construct $A_n=\int_0^\infty a_n(t)\,dt$.  For $|w|\le r<R$, finite
sums give
\[
 \left|\int_0^\infty f(t,z_0+w)\,dt-\sum_{n=0}^N A_nw^n\right|
 \le\frac{B(r/R)^{N+1}}{1-r/R},\qquad |A_n|\le BR^{-n}.
\]
To prove this, first work on $[\varepsilon,T]$, where everything is a
finite sum plus a uniform remainder.  Bound that remainder by the
geometric tail times $M(t)$.  The omitted end pieces have the same bound
with the corresponding tails of $M$.  Now make all three finite errors
small.  This proves termwise integration and constructs a local series
chart for the parameter function.  Its derivatives are the coefficient
computations from Chapter~\ref{ch:local-models}; no dominated-convergence
theorem is imported.

\section{Euler's gamma integral with uniform rational budgets}
For a rational complex $z$ with $\Re z>0$, define
\[
 \Gamma(z)=\int_0^\infty t^{z-1}E(-t)\,dt,
 \qquad t^{z-1}=E((z-1)L(t)).
\]
The same construction accepts computed inputs in a certified region
$a\le\Re z\le b$, where $a>0$ is rational.  On $0<t\le1$,
\[
 |t^{z-1}E(-t)|\le t^{a-1},\qquad e_0(\varepsilon)=\varepsilon^a/a.
\]
Choose an integer $m\ge\max(0,b-1)$.  On $t\ge1$, the absolute
integrand is bounded by $t^mE(-t)$.  Repeated finite integration by parts
on $[T,V]$, followed by the factorial bound for the exponential, gives
\[
 e_\infty(T)=E(-T)m!\sum_{j=0}^m\frac{T^j}{j!}.
\]
This expression tends to zero: for each fixed power, the exponential
series has a higher power as a positive term, giving a rational decay
bound, or repeated doubling of $T$ gives an effective search.  Computed
bounds such as $e_0,e_\infty$ are replaced by rational upper enclosures
when a numerical error budget is chosen.

On each $[\varepsilon,T]$, powers and exponentials have already proved
value, derivative, and continuity estimates, so the middle integral is
an ordinary rectangle computation.  Gamma is therefore constructed for
all these inputs, not introduced by assuming its improper integral exists.
Its standard normalization is recorded in
\href{https://dlmf.nist.gov/5.2}{DLMF, Section 5.2}.

\section{All parameter derivatives, by one envelope}
Fix $z_0$, and choose rational $0<\rho<a\le\Re z_0\le b$.
Expanding $E(wL(t))$ gives
\[
 a_n(t)=t^{z_0-1}E(-t)\frac{L(t)^n}{n!}.
\]
The positive exponential series proves
\[
 \frac{|L(t)|^n}{n!}\le\rho^{-n}
 \begin{cases}t^{-\rho},&0<t\le1,\\t^\rho,&t\ge1.\end{cases}
\]
Thus an integrable coefficient envelope is
\[
 M(t)=\begin{cases}t^{a-\rho-1}E(-t),&t\le1,\\
                   t^{b+\rho-1}E(-t),&t\ge1.
       \end{cases}
\]
Choose an integer $m\ge\max(0,b+\rho-1)$.  The preceding tail computations
apply to $M$, and $\int M\le B=1/(a-\rho)+m!$ is a rational bound.
The parameter-majorant argument supplies a gamma chart of radius $\rho$:
\[
 \Gamma^{(n)}(z_0)=\int_0^\infty t^{z_0-1}E(-t)L(t)^n\,dt,
 \qquad |\Gamma^{(n)}(z_0)|\le n!B\rho^{-n}.
\]
The integral on the right is a computed quantity with its own tails.
Every finite derivative order is now supported; no differentiation of an
uncontrolled infinite integral has occurred.

\section{Recurrence and continuation through known poles}
Integrate the derivative of $t^zE(-t)$ on a truncated interval.  The two
endpoint terms tend to zero by the same estimates, proving
\[
 \Gamma(z+1)=z\Gamma(z),\qquad \Gamma(1)=1.
\]
For a specific $z$ outside the nonpositive integers, choose $N$ so that
$\Re(z+N)>0$ and certify separation of each factor in
\[
 \boxed{\Gamma(z)=\frac{\Gamma(z+N)}{z(z+1)\cdots(z+N-1)}.}
\]
Any two allowable choices agree by the recurrence.  This is an actual
continuation algorithm, since shifting into the right half-plane supplies
new integral bounds, not just a recentering of an old disk.

At $z=-k$ the scalar value is replaced by its pole description.  Isolating
the factor $z+k$ in the denominator and using $\Gamma(1)=1$ gives residue
$(-1)^k/k!$.  The remaining analytic part can be computed by the local
series arithmetic.  For computed inputs, the required separation from
these poles must be supplied; no universal equality decision is asserted.
For positive real arguments gamma is positively separated from zero:
a lower rectangle on $[1,2]$ suffices.  Division by gamma elsewhere
requires its own separation proof; zero-freeness is not assumed here.

\section{Beta and a Gaussian normalization from the circle}
For positive real parameters $a,b$, define
\[
 B(a,b)=\int_0^1 t^{a-1}(1-t)^{b-1}\,dt.
\]
On $(0,1/2]$ the second factor has an explicit rational upper bound;
therefore the left tail is at most a constant times $\varepsilon^a/a$.
Reflection gives the right-tail bound.  The middle integral is bounded
quadrature.  These estimates work uniformly on separated parameter boxes.

\begin{theorem}[Beta--gamma comparison]
\label{thm:beta-gamma}
These independent computations satisfy
\[
 B(a,b)\Gamma(a+b)=\Gamma(a)\Gamma(b).
\]
\end{theorem}
\begin{proof}
Multiply the truncated gamma integrals and interpret the product as finite
rectangular sums in the positive quadrant.  Excluding small coordinate
strips and far tails changes it by the products of the one-dimensional
bounds already proved.  On each remaining bounded region the integrand
is bounded and uniformly continuous.  Triangulate and change the order of
integration; cells meeting the finitely many straight boundaries have
total area tending to zero with the mesh.  Thus the product is equivalently
computed by
\[
 \int_0^\infty E(-r)\int_0^r u^{a-1}(r-u)^{b-1}\,du\,dr.
\]
Here $r=u+v$ is just a shear with determinant $1$, verified on finite
parallelograms; deleting a diagonal tail $r>R$ is controlled by the union
of $u>R/2$ and $v>R/2$.  On each positive $r$ use the one-dimensional
substitution $u=rt$, first away from the two endpoints and then with
the beta endpoint bounds.  The inner integral is $r^{a+b-1}B(a,b)$.
Its remaining endpoint and infinity bounds are the gamma bounds for
$a+b$, proving the identity by finite error comparisons.
\end{proof}

At $a=b=1/2$, use $t=x^2$ and then the rational circle chart
$x=2u/(1+u^2)$, on truncated intervals.  Cancellation gives
\[
 B(1/2,1/2)=4\int_0^1\frac{du}{1+u^2}=\pi.
\]
The endpoint bounds justify both limiting substitutions.  Positivity and
the beta--gamma identity give $\Gamma(1/2)=\sqrt\pi$.  A final
substitution $t=x^2$ therefore proves
\[
 \int_0^\infty E(-x^2)\,dx=\frac{\sqrt\pi}{2}.
\]
The Gaussian normalization has been connected to the earlier circle
computation, without assuming a general two-dimensional change-of-variables
theorem.  Standard beta notation is in
\href{https://dlmf.nist.gov/5.12}{DLMF, Section 5.12}.

\section{The error function: convergent and asymptotic computations}
Define the entire local-series computation
\[
 \operatorname{erf}(z)=\frac2{\sqrt\pi}
       \sum_{n=0}^\infty\frac{(-1)^nz^{2n+1}}{n!(2n+1)}.
\]
On $|z|\le r$, consecutive absolute terms have ratio at most
$r^2/(n+1)$.  After that ratio is at most $1/2$, twice the first omitted
term bounds the tail.  Differentiated tails are obtained in the same way.
Finite integration of the exponential prefixes identifies the function
on the real axis as $2\int_0^x E(-t^2)dt/\sqrt\pi$.

For $x>0$, let $I_k(x)=\int_x^\infty E(-t^2)t^{-2k}dt$.
Finite integration by parts and the vanishing far endpoint give
\[
 I_k(x)=\frac{E(-x^2)}{2x^{2k+1}}-\frac{2k+1}{2}I_{k+1}(x).
\]
Repeated substitution proves, for every finite $N\ge1$,
\[
 I_0(x)=\frac{E(-x^2)}{2x}
   \sum_{k=0}^{N-1}\frac{(-1)^k(2k-1)!!}{(2x^2)^k}+R_N(x),
\]
where $(-1)!!=1$ and
\[
 0\le(-1)^NR_N(x)
 \le\frac{E(-x^2)}{2x}\frac{(2N-1)!!}{(2x^2)^N}.
\]
Indeed $R_N=(-1)^N(2N-1)!!\,I_N/2^N$, and the first integration-by-parts
inequality bounds $I_N$ by $E(-x^2)/(2x^{2N+1})$.

This gives the alternative computation
$1-\operatorname{erf}(x)=2I_0(x)/\sqrt\pi$.  It also illustrates an
important distinction: the asymptotic expression has a proved error for
each \emph{finite} $N$, but its terms eventually grow at fixed $x$.
It is not a convergent infinite series at that input.  Use it when its
finite bound meets the requested precision; otherwise use bounded
quadrature with tails or the convergent series.  The standard asymptotic
formula appears in \href{https://dlmf.nist.gov/7.12}{DLMF, Section 7.12}.

\chapter{Fourier Computations and the Basel Sum}\label{ch:fourier}

We need not begin with a space of arbitrary functions or a completeness
statement about an orthogonal basis.  A finite trigonometric sum is a
computation, and a positive finite kernel gives an explicit approximation
bound.  One worked periodic function will then identify the reciprocal-square
sum.

\section{Finite orthogonality}
Write $e_k(x)=E(i\pi kx)=C(kx)+iS(kx)$ for integer $k$.  The trigonometric
FTC gives
\[
 \frac12\int_{-1}^1 e_k(x)\,dx
 =\begin{cases}1,&k=0,\\0,&k\ne0.\end{cases}
\]
Products satisfy $e_ke_l=e_{k+l}$.  These identities determine integrals
of every finite trigonometric polynomial by finite linear algebra.

For a supplied period-$2$ continuous computation $f$, with a uniform
continuity certificate on the periodic interval, define
\[
 \widehat f(k)=\frac12\int_{-1}^1f(x)e_{-k}(x)\,dx.
\]
Each coefficient is an ordinary integral computation.  The word periodic
includes matching endpoint values and the continuity bound across that
join; this data is checked for the particular functions used here.

\section{A positive kernel with an explicit tail}
For $N\ge1$, define
\[
 K_N(t)=\frac1N\left|\sum_{j=0}^{N-1}e_j(t)\right|^2
       =\sum_{|k|<N}\left(1-\frac{|k|}{N}\right)e_k(t).
\]
The identity follows by counting pairs with a given difference.  Thus
$K_N\ge0$ and $\frac12\int_{-1}^1K_N=1$.
The finite geometric identity implies, for $0<|t|\le1$,
\[
 K_N(t)\le\frac1{N S(t/2)^2}\le\frac1{Nt^2}.
\]
For the last inequality, the midpoint identity makes $S$ concave on
$[0,1]$, where $S\ge0$: its centered second difference is
$2S(x)(C(h)-1)\le0$.  Rational Jensen inequalities on $[0,1/2]$ give
$S(t/2)\ge t$ for $0\le t\le1$.

The kernel computation
\[
 \sigma_N f(x)=\frac12\int_{-1}^1K_N(t)f(x-t)\,dt
\]
is also the finite trigonometric sum
\[
 \sum_{|k|<N}\left(1-\frac{|k|}{N}\right)\widehat f(k)e_k(x).
\]
This follows by expanding the finite kernel and shifting the integral;
periodicity makes any interval of length $2$ equivalent to $[-1,1]$.
The shift is justified by splitting at the finitely many period boundaries.

Let $|f|\le B$, and choose rational $0<\delta\le1$ so that periodic
arguments within $\delta$ have values within $\varepsilon$.  Split the
kernel integral at $|t|=\delta$.  Positivity and normalization bound the
near part by $\varepsilon$; the far part is at most
$2B/(N\delta^2)$.  Hence
\[
 \boxed{|\sigma_N f(x)-f(x)|\le\varepsilon+\frac{2B}{N\delta^2}.}
\]
This constructs a uniform trigonometric approximation with a terminating
precision schedule.  Computing the finitely many Fourier coefficients to
additional error $\eta/(2N)$ gives at most $\eta$ further output error.

\section{The periodic square}
Let $f(x)=x^2$ on $[-1,1]$, extended with period $2$.  Its endpoints match,
and it has periodic Lipschitz bound $2$, verified by splitting an increment
at the join if necessary.  Two integrations by parts give
\[
 \widehat f(0)=\frac13,\qquad
 \widehat f(k)=\frac{2(-1)^k}{\pi^2k^2}\quad(k\ne0).
\]
Here is the real calculation behind the coefficient: symmetry removes the
sine term, and with $a=\pi k$,
\[
 \int_0^1x^2\cos(ax)\,dx
 =-\frac2a\int_0^1x\sin(ax)\,dx
 =\frac{2(-1)^k}{a^2}.
\]
All boundary terms are explicit geometric sine and cosine values.

The series
\[
 F(x)=\frac13+\frac4{\pi^2}
          \sum_{k=1}^{\infty}\frac{(-1)^kC(kx)}{k^2}
\]
has a uniform tail at most $4/(\pi^2N)$ after term $N$, by the earlier
reciprocal-square tail.  It is therefore an independent continuous function
computation.  Its Fejer sums are the same finite sums as those of $f$:
termwise integration follows from the uniform tail, and finite
orthogonality gives the coefficients.  The kernel estimate shows that
these common sums converge to both $F$ and $f$.  Thus $F=f$.

\begin{theorem}[Basel identity]\label{thm:basel-geometric-overlap}
The reciprocal-square computation of Chapter~\ref{ch:infinite-series}
satisfies
\[
 \boxed{\sum_{k=1}^{\infty}\frac1{k^2}=\frac{\pi^2}{6}.}
\]
\end{theorem}
\begin{proof}
At the rational endpoint $x=1$, $C(k)=(-1)^k$ and $f(1)=1$.  The proved
series identity gives $1=1/3+(4/\pi^2)\sum_{k\ge1}1/k^2$.
\end{proof}

The chain of comparisons is now complete: rational circle areas define
$\pi$; convex secants compute sine and cosine derivatives; the FTC computes
Fourier coefficients; finite positive kernels identify a uniformly
convergent series; evaluating it identifies an independently constructed
number.  No Hilbert-space completeness or unrestricted Fourier convergence
theorem was needed.

\chapter{Corners, Jumps, and Delta Inputs}\label{ch:impulses}

A derivative computation need not always return a pointwise function.
At a corner, secant gaps need not shrink; at a jump, ordinary rational-point
derivatives miss an increment.  We retain those increments by allowing
finite atoms in addition to ordinary integrands.  The examples and their
finite evaluation rules come first; no space of all measures or all
distributions is required.

\section{A delta is an evaluation computation}
For a rational location $s$, define
\[
 \delta_s(\phi)=\phi(s).
\]
The input $\phi$ is a concrete test computation.  Continuous tests suffice
for a delta; a differentiated delta will require the corresponding
computed derivatives.  Compactly supported piecewise-polynomial tests,
with the necessary matching derivatives at joins, supply a first family.
They have finite bounds, moduli, and exact integrals.  For instance, on a
rational interval $[a,b]$, use $((t-a)(b-t))^{k+1}$ and set it to zero
outside: its first $k$ derivatives match at the joins.  Finite sums and
translates give explicit probes of any requested finite derivative order.

A finite signed collection of atoms is represented by a finite list
$(s_j,c_j)$ and evaluates a test as $\sum_jc_j\phi(s_j)$.  Adding an
ordinary integrable density $r$ gives the computation
\[
 T(\phi)=\int r(t)\phi(t)\,dt+\sum_jc_j\phi(s_j).
\]
Domain bounds and the integral construction for the displayed product are
part of this concrete description.

A computed location $s$ is also possible when the test has interval
continuity: evaluate it on the shrinking location enclosures as in
Chapter~\ref{ch:foundations}.  In particular, bisection of $x^2=2$ gives
$\delta_{\sqrt2}(\phi)$ without asking a rational sample to equal
$\sqrt2$.  This does not assert that a discontinuous step at an arbitrary
computed location can be evaluated at every rational input.  Such a
point evaluation requires a certified side or a separately chosen value
at the jump.

\section{Differentiating an explicitly joined function}
Suppose $y$ is supplied on finitely many pieces, with FTC and
integration-by-parts estimates for those branch computations.  Suppose its
one-sided values are computed and its jumps at the joins are
$c_j=y(s_j+)-y(s_j-)$.  For a compactly supported $C^1$ piecewise-polynomial
test $\phi$, integration by parts on each piece gives
\[
 -\int y(t)\phi'(t)\,dt
   =\int y'_{\mathrm{pieces}}(t)\phi(t)\,dt
      +\sum_jc_j\phi(s_j).
\]
Interior boundary terms are exactly the jumps; outer boundary terms vanish.
This proves the concrete derivative rule
\[
 \boxed{Dy=y'_{\mathrm{pieces}}+\sum_jc_j\delta_{s_j}.}
\]
Values assigned at the individual joins do not alter any of the integral
computations: cells covering finitely many points have arbitrarily small
total length and the functions have explicit local bounds.

The requirement on each piece is an FTC certificate, not just a pointwise
rational derivative.  Convex differences and the series constructions have
already supplied such certificates.  The ODE constructions in the next
chapter will supply further ones through integral equations.

\section{The ramp, the step, and the absolute value}
Write $H_s(t)=0$ for $t<s$ and $H_s(t)=1$ for $t>s$; choose a value at
$s$ only when point evaluation there is requested.  The join formula gives
\[
 DH_s=\delta_s.
\]
The continuous ramp $R_s(t)=\max(0,t-s)$ has ordinary branch derivatives
$0$ and $1$, with no value jump.  Hence
\[
 DR_s=H_s,\qquad D^2R_s=\delta_s,
\]
and
\[
 \boxed{D^2|t-s|=2\delta_s}
\]
by $|t-s|=2R_s(t)-(t-s)$.  The nonshrinking secant interval $[-1,1]$
for $|t|$ is not a delta value.  Its width $2$ records the slope jump,
which appears as an atom in the \emph{second} derivative.

For smooth enough concrete tests, define
\[
 (DT)(\phi)=-T(\phi'),\qquad
 (D^k\delta_s)(\phi)=(-1)^k\phi^{(k)}(s).
\]
This computes delta derivatives by ordinary test evaluation.  It does not
require defining products of singular objects.  In particular, no product
of two deltas is assigned a value here.

\section{The original hidden jump, now accounted for}
Let $\alpha$ be the positive root computation of $x^2=2$.  The two branch
joins in
\[
 F(x)=S(x)+\mathbf1_{\{x^2>2\}}
\]
have increments $-1$ at $-\alpha$ and $+1$ at $\alpha$.  Rational point
evaluation of its branches is decidable by the square comparison; the
one-sided values are computed by continuity of $S$.  The derivative is
therefore
\[
 DF=\pi C(x)\,dx+\delta_\alpha-\delta_{-\alpha}.
\]
On $[0,2]$ neither endpoint is an atom, and
\[
 \int_0^2 DF=\int_0^2\pi C(x)\,dx+1=1=F(2)-F(0).
\]
Here the notation for the integral means the ordinary integral plus the
explicit finite sum of atoms strictly inside the endpoints.  Endpoints at
atoms require a stated one-sided convention.  The missing increment has
become a finite part of the data rather than a failure of the theorem.

\section{Convex curvature measured by rational tents}
There is a second route to atoms which uses the secant construction directly.
For a rational polygonal test $\psi$ supported in $[a,b]$, with zero
endpoint values and knots $a=x_0<\cdots<x_N=b$, put
\[
 s_i=\frac{f(x_{i+1})-f(x_i)}{x_{i+1}-x_i},\qquad
 \mu_f(\psi)=\sum_{i=1}^{N-1}\psi(x_i)(s_i-s_{i-1}).
\]
Here $f$ is any supplied convex computation on a larger interval; it need
not have a pointwise derivative.  Rearrangement gives
\[
 \mu_f(\psi)=\sum_i f(x_i)
          (\text{slope of $\psi$ to the right}
           -\text{slope of $\psi$ to the left}),
\]
including the endpoint jumps from and to zero slope.  Thus inserting a knot
where $\psi$ is linear changes nothing.  The definition is independent of
its polygonal description.

Convexity gives $\mu_f(\psi)\ge0$ for $\psi\ge0$.  Outer secants give a
rational bound $V$ for the total increase of the $s_i$, and hence
\[
 |\mu_f(\psi)|\le V\|\psi\|_\infty.
\]
Approximate a continuous test by rational polygonal tests with uniform
error $2^{-n}$ and evaluate the finite expression to error $2^{-n}$.
Enlarging by $(V+1)2^{-n}$ and taking finite intersections constructs
$\mu_f(\phi)$.  Cross-overlap follows from the displayed norm bound.
This is the curvature-measure computation associated with this particular
convex function, not an abstract measure-existence assertion.

For the tent of height $1$ with vertices $u<v<w$,
\[
 \mu_f(\tau)=\frac{f(w)-f(v)}{w-v}
                  -\frac{f(v)-f(u)}{v-u}.
\]
In the symmetric case this is precisely the gap between the left and right
derivative secants.  For $f(t)=|t-s|$ it is $2\tau(s)$; consequently
$\mu_f=2\delta_s$ on all the test computations just constructed.

\section{A finite second-order FTC}
For rational $a<x<b$, let
\[
 G_{a,b}(x,t)=
 \begin{cases}
 (t-a)(b-x)/(b-a),&a\le t\le x,\\
 (x-a)(b-t)/(b-a),&x\le t\le b,\\
 0,&\text{otherwise}.
 \end{cases}
\]
This is a rational tent multiple as a function of $t$.  Substitution in
the preceding finite formula gives
\[
 \boxed{f(x)=\frac{b-x}{b-a}f(a)+\frac{x-a}{b-a}f(b)
                -\mu_f(G_{a,b}(x,\cdot)).}
\]
Thus the endpoint values and the curvature computation reconstruct $f$.
The affine part is explicit.  At a corner this formula continues to work
without a scalar derivative at the corner.

For a compactly supported, globally $C^2$ piecewise-polynomial test,
finite summation by parts also gives
$\mu_f(\phi)=\int f\phi''$.  Indeed the polygonal interpolants produce
centered second differences of the test; the twice-integrated increment
bound for the test compares these with $h\phi''$ uniformly.  The bounded
function $f$ contributes at most its bound times the total error.
This identifies the curvature computation as $D^2f$ on these tests.
The same proof applies to any supplied test with that twice-integrated
increment identity and a uniform modulus for its second derivative;
pointwise rational differentiability alone is not the required data.

\section{Narrow pulses give an equivalent computation}
For rational $\varepsilon>0$ let
$p_{\varepsilon,s}=\varepsilon^{-1}\mathbf1_{[s,s+\varepsilon]}$.
A test with Lipschitz bound $L$ satisfies
\[
 \left|\int p_{\varepsilon,s}(t)\phi(t)\,dt-\phi(s)\right|
       \le L\varepsilon/2.
\]
The proof integrates the bound $L|t-s|$ on the pulse interval.  This is an
explicit equivalence between two ways to evaluate a delta: point evaluation
and narrowing ordinary rectangles.  It is convergence on tests, not
pointwise convergence to an infinitely high function.

\chapter{Linear ODEs and the Fundamental Matrix}
\label{ch:differential-equations}

A fundamental matrix is a reusable computation depending only on the
coefficient matrix.  It propagates initial states, accumulates ordinary
forcing, and propagates delta inputs.  We construct it from iterated finite
integrals and factorial error bounds.

\section{Coefficient data and finite norms}
Fix a rational interval $[a,b]$.  The entries of $A(t)$ are supplied
rational-input computations with interval-continuity certificates and a
rational bound $\|A(t)\|\le M$.  Use the maximum norm for vectors and the
maximum absolute row-sum norm for matrices; finite algebra proves
$\|AB\|\le\|A\|\|B\|$ and $\|Av\|\le\|A\|\|v\|$.

Polynomial matrices are the first examples.  Separated rational functions
and the elementary computations of the preceding chapters provide others.
A finite set of coefficient breakpoints is handled by solving on each
piece and composing the resulting evolutions.

\section{Peano--Baker prefixes}
For rational $s,t\in[a,b]$, define
\[
 P_0(t,s)=I,\qquad P_{n+1}(t,s)=\int_s^t A(u)P_n(u,s)\,du,
 \qquad \Phi_N(t,s)=\sum_{n=0}^NP_n(t,s).
\]
For fixed $n$ these are finite nested integral computations.  Their
continuity data is obtained recursively from products and the
bounded-integrand endpoint estimate.  If $A$ is a rational-polynomial
matrix, all $P_n$ are computed by finite polynomial integration.
Matrix multiplication is kept in the displayed order; matrices at distinct
times are not assumed to commute.

For $h=|t-s|$, induction using polynomial integration gives
\[
 \|P_n(t,s)\|\le\frac{(Mh)^n}{n!}.
\]
The argument works with oriented integrals in either time direction.
If $z=M(b-a)$ and $N+2\ge2z$, every finite tail after $N$ is bounded by
\[
 \boxed{2\frac{z^{N+1}}{(N+1)!}.}
\]
Compute the finite prefix to an additional rational error budget and
stabilize the resulting boxes.  This constructs the matrix computation
$\Phi(t,s)$ uniformly on the interval.  A rational upper bound $B$ for
$E(M(b-a))$ bounds every $\|\Phi(t,s)\|$.

\begin{figure}[htbp]\centering
\peanoBakerWordsAnimation
\caption{Ordered matrix products and their finite integral prefixes.}
\end{figure}

The finite recursion, with its uniform tails, gives
\[
 \Phi(t,s)=I+\int_s^t A(u)\Phi(u,s)\,du.
\]
The accumulation estimate therefore constructs its derivative:
\[
 \partial_t\Phi(t,s)=A(t)\Phi(t,s),\qquad\Phi(s,s)=I.
\]
No general differentiation-of-a-limit assertion has been used.

\section{Uniqueness belongs to the integral construction}
Suppose a bounded computed vector $w$ satisfies
\[
 w(t)=\int_s^t A(u)w(u)\,du.
\]
Starting with any rational bound $\|w\|\le C$, repeated substitution gives
\[
 \|w(t)\|\le C\frac{(M|t-s|)^n}{n!}
\]
for every integer $n$.  The right side tends to zero with an explicit
factorial bound, so $w$ is equivalent to zero.

For the solutions considered here, the integral identity is proved by
construction.  We do not infer it from a bare equation at every rational
point; a function hidden across a cut could satisfy such a pointwise
equation and fail its integral identity.

The same uniqueness argument gives, for all rational $s,u,t$,
\[
 \boxed{\Phi(t,u)\Phi(u,s)=\Phi(t,s).}
\]
Both sides have the same integral equation and initial value when regarded
as functions of $t$.  The identity follows for every starting location by
subtracting their integral equations there.  In particular
$\Phi(s,t)$ is the inverse of $\Phi(t,s)$; no determinant inversion is
needed to construct it.

The endpoint estimate gives
$\|\Phi(u,s)-I\|\le MB|u-s|$.  Together with composition it yields
\[
 \|\Phi(t,s)-\Phi(t,u)\|\le MB^2|s-u|.
\]
This supplies continuity in the source time as well as the target time,
and permits computed times with domain enclosures.

\section{Ordinary forcing and variation of constants}
For a supplied interval-continuous forcing $r$, put
\[
 y(t)=\Phi(t,a)y_a+\int_a^t\Phi(t,s)r(s)\,ds.
\]
All entries of this integrand have the continuity and boundedness data just
proved.  The finite integral recursions show
\[
 y(t)=y_a+\int_a^t\bigl(A(u)y(u)+r(u)\bigr)\,du.
\]
One verifies the double-integral rearrangement on finite subdivisions of
the triangle $a\le s\le u\le t$.  The cells meeting its diagonal have
total area bounded by a constant times $(b-a)\mesh(P)$; the integrands
have rational bounds, and the other cells are finite iterated rectangles.
This proves the needed interchange without a general measure-theoretic
Fubini theorem.

Accumulation gives $Dy=Ay+r$, and the preceding uniqueness estimate proves
uniqueness among the certified integral solutions.  Thus the displayed
formula is the solution computation, not a formula justified by an
unconstructed classical solution.

The same estimate controls numerical replacements.  If a proposed
computation $\widetilde y$ has a certified integral identity with initial
error at most $e$ and additional forcing (residual) bounded by $\rho$,
then variation of constants gives
\[
 \|\widetilde y(t)-y(t)\|\le B\bigl(e+(b-a)\rho\bigr).
\]
A merely sampled residual would not suffice: its bound must hold on the
cells, with the integral identity certified for that computation.

\section{Delta inputs are finite state updates}
For a finite list of impulse times $a<\tau_j<b$ and vectors $c_j$, solve
\[
 Dy=A(t)y+r(t)+\sum_jc_j\delta_{\tau_j}.
\]
Initially take the times rational.  For times away from the impulses the
solution is
\[
 \boxed{y(t)=\Phi(t,a)y_a+\int_a^t\Phi(t,s)r(s)\,ds
                  +\sum_{a<\tau_j<t}\Phi(t,\tau_j)c_j.}
\]
Its one-sided states satisfy
$y(\tau_j+)=y(\tau_j-)+c_j$ because $\Phi(\tau_j,\tau_j)=I$.
The piecewise integration-by-parts rule of Chapter~\ref{ch:impulses}
proves exactly the displayed equation on tests.  The ordinary branches
have the needed integration identities from their integral equations and
the corresponding finite product rule with a smooth test.

One computation evaluates the formula directly.  Another evolves to each
impulse, adds its vector, and continues.  The composition law proves their
equivalence.  A computed nonrational impulse location can also be used
when ordering and separated evaluation times are supplied; one-sided
states are retained instead of forcing an undecidable comparison at the
location itself.

Define the causal matrix kernel
\[
 G(t,s)=\begin{cases}0,&t<s,\\\Phi(t,s),&t>s.\end{cases}
\]
Its jump is $I$, and hence
\[
 \boxed{(D_t-A(t))G(t,s)=I\delta_s.}
\]
Each column is the response to a unit impulse in one state coordinate.
This identity is the reason for including delta inputs with the fundamental
matrix.

For a rectangular pulse of width $\varepsilon$, after the pulse ends the
response is $\varepsilon^{-1}\int_s^{s+\varepsilon}\Phi(t,u)c\,du$.
The source-time bound gives its difference from $\Phi(t,s)c$ as at most
$MB^2\|c\|\varepsilon/2$.  This is agreement between the state-update
construction and narrowing ordinary forcing, not uniform convergence
through a discontinuity.

\section{Three explicit kernels}
\subsection{The free second derivative}
For the state $(y,y')$ and matrix
$A=\left(\begin{smallmatrix}0&1\\0&0\end{smallmatrix}\right)$,
\[
 \Phi(t,s)=\begin{pmatrix}1&t-s\\0&1\end{pmatrix}.
\]
An impulse $(0,1)^{\mathsf T}$ gives $y=(t-s)_+$, so $D^2y=\delta_s$.
The function $|t-s|$ solves the equation with $2\delta_s$ but differs
from its zero-past solution by an affine function.  Initial data decides
that homogeneous ambiguity.

\subsection{The oscillator}
For $y''+\pi^2y=0$,
\[
 \Phi(t,s)=
 \begin{pmatrix}
 C(t-s)&S(t-s)/\pi\\
 -\pi S(t-s)&C(t-s)
 \end{pmatrix}.
\]
Its entries, derivatives, and composition are already certified by the
circle and FTC chapters.  The impulse in the velocity coordinate gives
\[
 g_s(t)=\begin{cases}0,&t<s,\\S(t-s)/\pi,&t>s,\end{cases}
 \qquad (D^2+\pi^2)g_s=\delta_s.
\]
The position is continuous and the velocity jumps by $1$.  A position jump
would instead introduce a derivative of delta into the second derivative.

\subsection{A variable coefficient: the Airy equation}
For $y''-ty=0$, substitution of $y=\sum a_nt^n$ gives
\[
 a_2=0,\qquad a_{n+3}=\frac{a_n}{(n+3)(n+2)}.
\]
Take $u(0)=1,u'(0)=0$ and $v(0)=0,v'(0)=1$.  Thus
\[
 u(t)=1+t^3/6+t^6/180+\cdots,\qquad
 v(t)=t+t^4/12+t^7/504+\cdots.
\]
On $|t|\le R$, each of the three coefficient strands eventually has
successive absolute-term ratio at most $1/2$, by taking
$(n+3)(n+2)\ge2R^3$.  This gives geometric tail bounds; differentiated
strands have the same property after increasing the index.  Thus $u,v$
and their derivatives are actual series computations on every bounded
interval.

Coefficient multiplication proves $uv'-u'v=1$: its formal derivative is
$uv''-u''v=0$, so all nonconstant coefficients vanish, and its constant
coefficient is $1$.  Tails justify the identity for the computations.
Consequently
\[
 g_s(t)=\begin{cases}
 0,&t<s,\\
 u(s)v(t)-v(s)u(t),&t>s
 \end{cases}
\]
is continuous at $s$ and its first derivative jumps by $1$.  Hence
\[
 (D^2-t)g_s=\delta_s.
\]
This constructs a variable-coefficient impulse response by rational
recurrences and finite determinant algebra.  It is equivalent to the
Peano--Baker computation by integral-equation uniqueness.

\section{The FTC as the zero-coefficient case}
When $A=0$, the fundamental matrix is $I$.  Ordinary forcing and the
finite atoms give
\[
 y(t)-y(a)=\int_a^t r(s)\,ds+\sum_{a<\tau_j<t}c_j,
\]
with endpoints away from impulses.  Thus integration of a derivative is
inversion of $D$, up to a constant; linear ODE evolution is inversion of
$D-A$, up to a homogeneous solution.  Both are identities between explicit
computations with initial data, finite updates, and shrinking errors.

\chapter{Special Functions from Equations and Integrals}
\label{ch:special-equations}

A differential equation is useful only after its chosen solution has been
specified and its computation bounded.  A recurrence supplies local data;
the fundamental matrix moves that data along certified paths.  Bessel
and hypergeometric functions illustrate the two parts, and a Bessel zero
shows how a special-function value can itself become an inverse computation.

\section{Bessel functions of integer order}
For an integer $m\ge0$, define
\[
 J_m(z)=\sum_{k\ge0}\frac{(-1)^k(z/2)^{2k+m}}{k!(k+m)!}.
\]
Every coefficient is rational.  On $|z|\le R$, the absolute terms have
successive ratio
\[
 \frac{R^2}{4(k+1)(k+m+1)}.
\]
This decreases to zero; after it is at most $1/2$, the tail is at most
twice the first omitted majorant term.  Multiplying by any fixed falling
factorial gives a derivative-tail sequence with the same eventual property.
Consequently these are entire function and derivative computations.  To
obtain a geometric coefficient chart of any chosen radius, evaluate the
absolute series majorant at a strictly larger radius and bound each
coefficient by that weighted sum.

Coefficient comparison proves
\[
 z^2J_m''+zJ_m'+(z^2-m^2)J_m=0,
 \qquad J_0'=-J_1.
\]
For example the coefficient of degree $2k+m$ in the equation is
$4k(k+m)a_k+a_{k-1}=0$ for $k\ge1$; the lowest coefficient is free
and is set to $1/(2^m m!)$.  This both verifies the equation and records
which local solution was selected.  Standard notation and the equation
are in \href{https://dlmf.nist.gov/10.2}{DLMF, Section 10.2}.

On a path avoiding zero, the state $(J_m,J_m')$ satisfies
\[
 Y'=\begin{pmatrix}0&1\\m^2/z^2-1&-1/z\end{pmatrix}Y.
\]
Start at a nonzero point with values from the entire series.  Pull back
each rational polygonal path segment $z=p+t(q-p)$ to $0\le t\le1$:
its coefficient matrix is $(q-p)A(p+t(q-p))$.  A certified positive
separation from zero bounds all entries and supplies their continuity.
The Peano--Baker construction then propagates the state.  Integral-equation
uniqueness proves agreement with the series.  It is unnecessary, and
incorrect, to apply the bounded-coefficient evolution theorem at $z=0$;
the series handled that point separately.

\section{A second Bessel computation}
The circle functions give the bounded integral
\[
 B(z)=\int_0^1 E(izC(t))\,dt.
\]
The factorial majorant is uniform for $0\le t\le1$ on every bounded
$z$-disk.  Thus it can be integrated term by term.  Symmetry cancels odd
powers, and finite trigonometric integration gives
\[
 \int_0^1 C(t)^{2k}\,dt=\frac{\binom{2k}{k}}{4^k}.
\]
For completeness, differentiating $S(t)C(t)^{2k-1}$ and integrating its
zero endpoint difference yields the recurrence
$M_k=(2k-1)M_{k-1}/(2k)$, with $M_0=1$.
The product has the uniform increment estimates already proved for sine
and cosine, so this use of the FTC is justified.
Substitution into the exponential coefficients shows $B=J_0$.
For real $z$ its imaginary part cancels by $t\mapsto1-t$; hence
\[
 J_0(z)=\int_0^1\cos(zC(t))\,dt,
\]
where $\cos w=(E(iw)+E(-iw))/2$.  This is the normalized Bessel integral
of \href{https://dlmf.nist.gov/10.9}{DLMF, Section 10.9}, proved here
by finite moments and a uniform tail.

\section{Computing a zero without a general root theorem}
\label{sec:bessel-zero}
Exact alternating-series bounds give
\[
 J_0(2)>\frac29>0,\qquad J_0(3)<-\frac{4199}{16384}<0.
\]
The first bound uses terms through $k=3$; the second uses terms through
$k=4$.  Although the first absolute term at $3$ need not decrease, the
remaining tail after these cutoffs does decrease and has the indicated
sign.  At $0\le x\le3$, writing $q=x^2$, the tail in $J_1$ after its
four displayed terms gives
\[
 J_1(x)\ge\frac x2
       \left(1-\frac q8+\frac{q^2}{192}-\frac{q^3}{9216}\right).
\]
The polynomial in parentheses decreases on $[0,9]$: its derivative is
$-1/8+q/96-q^2/3072<0$.  Its value at $9$ is $223/1024$.
Therefore $J_1(x)\ge223/1024$ on $[2,3]$, and the certified FTC gives
\[
 J_0(u)-J_0(v)\ge\frac{223}{1024}(v-u)
 \quad(2\le u<v\le3).
\]
The endpoint signs and this separation make the trisection inverse
construction of Chapter~\ref{ch:foundations} terminate.  They compute a
unique zero in $[2,3]$ as nested rational intervals.  Also
$J_0(x)\ge1-x^2/4>0$ for $0\le x<2$, so this is its first positive
zero.  The root is a new number computation, not a symbolic placeholder
chosen by a completeness theorem.

\section{A recurrence with a finite convergence disk}
For rational parameters $a,b,c$ with $c\notin\{0,-1,-2,\ldots\}$, put
\[
 h_0=1,\qquad
 (n+1)(n+c)h_{n+1}=(n+a)(n+b)h_n,
 \qquad F(z)=\sum_{n\ge0}h_nz^n.
\]
If an upper parameter makes a numerator zero, the recurrence terminates
and gives a polynomial.  Otherwise, for a chosen rational $r<1$, bound
the term ratio by
\[
 r\frac{(n+|a|)(n+|b|)}{(n+1)(n-|c|)}\qquad(n>|c|).
\]
Choose rational $r<\theta<1$.  After clearing the positive denominator,
the assertion that this ratio is at most $\theta$ is a quadratic
polynomial inequality in $n$ with positive leading coefficient.  Choose
an integer $N$ large enough that its expansion in $n=N+j$ has nonnegative
coefficients.  That finite check proves the ratio bound for every $j\ge0$.
The tail after $N$ is then bounded by its first omitted majorant term
divided by $1-\theta$.  Derivative tails are handled by the same method.

To supply a chart on $|z|<R<1$, apply this ratio argument at $R$ and
bound the finitely many earlier weighted coefficients.  This gives a
rational $M$ with $|h_n|\le M R^{-n}$ for all $n$.
Coefficient comparison now proves
\[
 \boxed{z(1-z)F''+[c-(a+b+1)z]F'-abF=0.}
\]
The recurrence and $h_0=1$ identify the local solution, denoted
${}_2F_1(a,b;c;z)$; see
\href{https://dlmf.nist.gov/15.2}{DLMF, Section 15.2} and
\href{https://dlmf.nist.gov/15.10}{Section 15.10} for the standard notation.

\section{Continuation uses the equation, not repeated recentering alone}
Away from $0$ and $1$ the preceding equation is a first-order system with
\[
 A(z)=\begin{pmatrix}
 0&1\\
 ab/[z(1-z)]&((a+b+1)z-c)/[z(1-z)]
 \end{pmatrix}.
\]
The convergent series gives starting data at a nonzero point inside the
unit disk.  A finite path with certified separation from $0,1$ gives
bounded coefficient computations on each pulled-back segment; the
fundamental matrix propagates the state along the path.

It also supplies fresh analytic charts.  Here is a finite majorant
argument for any rational-coefficient system at an ordinary point.
Suppose locally $A(z_0+w)=\sum A_jw^j$ with
$\|A_j\|\le M R^{-j}$, obtained by polynomial and reciprocal charts.
The solution coefficients satisfy
\[
 (n+1)y_{n+1}=\sum_{j=0}^n A_jy_{n-j}.
\]
Choose $0<r\le R/2$ with $2Mr\le1$, and a rational
$C\ge\|y_0\|$.  Induction gives
\[
 \|y_n\|\le Cr^{-n}:
 \quad \|y_{n+1}\|\le\frac{2MC}{n+1}r^{-n}
                                      \le Cr^{-n-1}.
\]
On $|w|\le r/2$ the local series and all requested derivative tails
are therefore effective.  Its integral equation identifies it with the
propagated state.  Coefficient separation from the singularities allows
this bound to be renewed at successive centers.

Distinct paths may produce distinct branches.  The finite overlap and
path-deformation certificates of Chapter~\ref{ch:complex-paths} say when
they agree; there is no assumption of global single-valuedness.

\section{The beta integral identifies the same local function}
For positive rational $b,c-b$ and rational $a$, define on $|z|<1$
\[
 H(z)=\frac1{B(b,c-b)}\int_0^1
       t^{b-1}(1-t)^{c-b-1}(1-zt)^{-a}\,dt.
\]
Use the local logarithm to define the last factor.  Its coefficient
recurrence is $(n+1)d_{n+1}=(n+a)d_n$, $d_0=1$, obtained from
$(1-w)D(1-w)^{-a}=a(1-w)^{-a}$.  The ratio estimate just used supplies
uniform tails on smaller disks.  The positive beta weight provides the
endpoint majorant, so the parameter-integral theorem applies.

The beta--gamma identity and gamma recurrence give
\[
 \frac{B(b+n,c-b)}{B(b,c-b)}
       =\frac{b(b+1)\cdots(b+n-1)}{c(c+1)\cdots(c+n-1)}.
\]
Thus integrating the binomial coefficients yields exactly the $h_n$
above, and $H=F$.  This proves equivalence between an improper integral,
a convergent series, and a propagated ODE computation.
In particular the elliptic integral developed later has
\[
 K(m)=\frac\pi2\,{}_2F_1(1/2,1/2;1;m).
\]
The later geometric substitution gives its independent bounded quadrature;
the coefficient identity here explains why its equation has the displayed
hypergeometric form.

\chapter{Transforms, Oscillatory Tails, and Zeta}
\label{ch:transforms-zeta}

The same function can be computed by a transform, an equation, or a finite
sum with endpoint corrections.  Three examples show why the foundation
must keep track of the error appropriate to each operation.  None requires
an ambient space of arbitrary continuous or square-integrable functions.

\section{Laplace transforms with an explicit half-plane}
Suppose a concrete, locally integrable computation $f$ has a rational
exponential bound $|f(t)|\le B E(\alpha t)$ on $t\ge0$.  On a certified
half-plane $\Re s\ge\alpha+\sigma$, with rational $\sigma>0$, define
\[
 \mathcal Lf(s)=\int_0^\infty E(-st)f(t)\,dt.
\]
Its far tail is at most $B E(-\sigma T)/\sigma$.  The middle integral
has the data of the particular $f$.  This constructs the transform.
Multiplication of the integrand by $t^k$ still gives a computable
exponential tail, by the polynomial-exponential calculation in
Chapter~\ref{ch:improper-parameters}.  Expanding $E(-ht)$ with a smaller
half-plane margin gives the parameter-majorant proof of
\[
 (\mathcal Lf)^{(k)}(s)=(-1)^k\int_0^\infty t^kE(-st)f(t)\,dt.
\]
The half-plane is part of the certificate, not an unspecified region of
convergence.

If $f$ is a certified piecewise primitive with finitely many jumps and
its ordinary derivative has the needed tails, finite integration by parts
and the explicit jump contributions give
\[
 \mathcal L(Df)(s)=s\mathcal Lf(s)-f(0+),
\]
where the transform includes interior impulse masses and no impulse at
the initial endpoint.  Extending $f$ by zero to negative times instead
adds $f(0+)\delta_0$ to its full-line derivative.  This distinguishes
the two initial-value conventions rather than assigning a hidden half
weight to the delta at zero.

For an interior impulse $\delta_\tau$, its transform is $E(-s\tau)$.
In particular, ordinary integrations by parts give
\[
 \mathcal L(S)(s)=\frac\pi{s^2+\pi^2},\qquad
 \mathcal L(C)(s)=\frac{s}{s^2+\pi^2}\quad(\Re s>0).
\]
The denominators are separated on any supplied closed box inside that
half-plane.  These formulas compare the oscillator's impulse kernel with
a rational computation; they are not applications of an unproved transform
inversion theorem.

For a constant matrix $A$, the zero-past impulse response is $E(tA)$
for $t>0$.  On $\Re s>\|A\|$, entrywise integration gives
\[
 (sI-A)\mathcal L(E(\cdot A))(s)=I.
\]
The finite geometric matrix series
$s^{-1}\sum_{n\ge0}(A/s)^n$, with its norm tail, independently constructs
$(sI-A)^{-1}$.  Uniqueness of that inverse proves their equivalence.  Thus
the familiar resolvent is a second computation of the fundamental matrix's
Laplace transform.

\section{An oscillatory integral with no absolute-integral bound}
Extend $S(t)/t$ at zero by the series value $\pi$.  For $0<T<U$,
integration by parts using $DS=\pi C$, $DC=-\pi S$ gives
\[
 \left|\int_T^U\frac{S(t)}t\,dt\right|\le\frac{2}{\pi T}.
\]
Indeed the endpoint terms have bounds $1/(\pi T)$ and $1/(\pi U)$,
and the remaining cosine integral is bounded by
$(1/T-1/U)/\pi$.  Therefore bounded quadrature and this cancellation
bound construct $I=\int_0^\infty S(t)/t\,dt$.

To identify it, first use a positive real damping parameter $a$:
\[
 F(a)=\int_0^\infty E(-at)\frac{S(t)}t\,dt.
\]
On $a\ge a_0>0$ the exponential tails justify parameter differentiation,
giving $F'(a)=-\pi/(a^2+\pi^2)$ by the Laplace calculation above.
Also $|F(a)|\le\pi/a$ from the chord bound $|S(t)|\le\pi t$.
Integrating this derivative between positive rational endpoints and sending
the larger one to infinity with that explicit bound yields
\[
 F(a)=\int_a^\infty\frac\pi{u^2+\pi^2}\,du
     =\frac\pi2-A(a/\pi)
\]
when $0<a\le\pi$; scaling and the previously computed angle integral
prove the last equality.

The damping can be removed with a finite estimate.  The function
$E(-at)/t$ is positive decreasing; the same integration-by-parts argument
bounds its oscillatory tail uniformly for all $a\ge0$ by $2/(\pi T)$.
On $[0,T]$, the bound $|1-E(-at)|\le at$ gives
\[
 |F(a)-I|\le aT+\frac4{\pi T}.
\]
First choose a rational $T$, then a sufficiently small positive rational
$a$.  This proves $F(a)\to I$ without an absolute domination on the
infinite interval.  Since $A(a/\pi)\to0$ by its continuity estimate,
\[
 \boxed{\int_0^\infty\frac{S(t)}t\,dt=\frac\pi2.}
\]
For the normalized sinc $S(t)/(\pi t)$, the positive-half-line integral
is $1/2$ and the whole-line integral is $1$.  The endpoint and normalization
are stated explicitly.

\section{A Fourier transform proved by an ODE}
Let $f(x)=E(-\pi x^2)$ and define for real $\xi$
\[
 H(\xi)=\int_{-\infty}^\infty f(x)E(-2\pi i\xi x)\,dx.
\]
The gamma normalization and scaling give $H(0)=1$.  Each parameter
derivative is controlled by a polynomial times the Gaussian.  For complex
increments of size at most $r$, the envelope gains at most $E(2\pi r|x|)$.
Once $|x|\ge4r$, this is still bounded together with the Gaussian by
$E(-\pi x^2/2)$.  The remaining bounded interval is handled separately.
This provides explicit integrable majorants for the parameter series.

Differentiate under the integral by those bounds, and integrate the
identity $f'(x)=-2\pi xf(x)$ by parts on finite symmetric intervals.
Its boundary terms tend to zero with their Gaussian bounds.  The result is
\[
 H'(\xi)=-2\pi\xi H(\xi),\qquad H(0)=1.
\]
The integral-equation uniqueness of the ODE chapter identifies its
computation as
\[
 \boxed{\int_{-\infty}^\infty E(-\pi x^2)E(-2\pi i\xi x)\,dx
                         =E(-\pi\xi^2).}
\]
Multiplying the integrand by $x^k$ gives
$(-2\pi i)^{-k}$ times the $k$th derivative on the right.  Thus each
Gaussian times a specified polynomial has its transform computed by finite
differentiation and polynomial algebra.  This is a genuine Fourier result
without constructing a space of arbitrary $L^2$ functions.

\section{Bernoulli polynomials are finite rational recurrences}
To go beyond a series' original convergence region we shall use finite
summation by parts.  Define polynomials $B_n(x)$ by
\[
 B_0=1,\qquad B_n'=nB_{n-1},\qquad\int_0^1 B_n(x)\,dx=0\quad(n\ge1).
\]
Polynomial integration fixes the constant term uniquely, so every coefficient
is rational.  The first polynomials are
\[
 B_1=x-\tfrac12,\quad B_2=x^2-x+\tfrac16,
 \quad B_4=x^4-2x^3+x^2-\tfrac1{30}.
\]
Induction using the recurrence and the zero mean proves
$B_n(1-x)=(-1)^nB_n(x)$.  Also $B_n(1)=B_n(0)$ for $n\ge2$.
Consequently the odd endpoint values vanish for odd $n>1$.
Put $b_{2k}=B_{2k}(0)$ and $\widetilde B_n(t)=B_n(t-\lfloor t\rfloor)$.
For integration this is simply a separate polynomial on each unit cell;
no discontinuous floor computation at a computed input is needed.

\section{Euler--Maclaurin as a finite identity}
\label{thm:finite-euler-maclaurin}
For integers $N<M$ and a specific function $f$ with certified derivatives
and integration-by-parts identities through order $2p$, one has
\begin{align*}
 \sum_{n=N}^{M-1} f(n)
 &=\int_N^M f(t)\,dt+\frac{f(N)-f(M)}2\\
 &\quad+\sum_{k=1}^p\frac{b_{2k}}{(2k)!}
                \bigl(f^{(2k-1)}(M)-f^{(2k-1)}(N)\bigr)\\
 &\quad-\frac1{(2p)!}\int_N^M\widetilde B_{2p}(t)f^{(2p)}(t)\,dt.
\end{align*}
\begin{proof}
On $[n,n+1]$, integration by parts gives
\[
 \int_n^{n+1}B_1(t-n)f'(t)\,dt
       =\tfrac12(f(n)+f(n+1))-\int_n^{n+1}f(t)\,dt.
\]
Sum these identities.  Repeatedly use
$\widetilde B_j=(\widetilde B_{j+1}/(j+1))'$ inside each cell to integrate
the remaining term by parts.  Matching endpoint values cancel the internal
boundaries for $j+1\ge2$, and the odd endpoint values above eliminate the
odd corrections.  After $2p$ derivatives the displayed remainder is what
remains.  Every step is a finite polynomial-weighted integral identity.
\end{proof}

This is the finite formula underlying the usual
\href{https://dlmf.nist.gov/2.10}{Euler--Maclaurin expansion}.  Its remainder
is part of the computation, not an unspecified asymptotic error.

\section{Zeta beyond the defining sum}
On a certified half-plane $\Re s\ge\sigma>1$, the series
\[
 \zeta(s)=\sum_{n\ge1}E(-sL(n))
\]
has tail at most $N^{1-\sigma}/(\sigma-1)$ after $N$ terms, by finite
comparison with the integral of $t^{-\sigma}$.  This first constructs
zeta only where the series has that bound.

Apply the finite Euler--Maclaurin identity to $f(t)=t^{-s}$, whose
uniform increment and derivative identities follow from the elementary
exponential and logarithm computations.  Here
\[
 f^{(j)}(t)=(-1)^j s^{\overline j}t^{-s-j},\qquad
 s^{\overline j}=s(s+1)\cdots(s+j-1)
\]
is a rising factorial, with $s^{\overline0}=1$.
Sending the upper endpoint to infinity first in $\Re s>1$ gives
\begin{align*}
 \zeta(s)&=\sum_{n=1}^{N-1}n^{-s}
       +\frac{N^{1-s}}{s-1}+\frac12N^{-s}\\
 &\quad+\sum_{k=1}^p\frac{b_{2k}}{(2k)!}
                   s^{\overline{2k-1}}N^{-s-2k+1}
       +R_{N,p}(s),\\
 R_{N,p}(s)&=-\frac{s^{\overline{2p}}}{(2p)!}
              \int_N^\infty\widetilde B_{2p}(t)t^{-s-2p}\,dt.
\end{align*}
Let $C_{2p}$ be the sum of the absolute coefficients of $B_{2p}$, a
rational upper bound on its absolute value in $[0,1]$.  For
$\Re s\ge\sigma>1-2p$, the remainder satisfies
\[
 \boxed{|R_{N,p}(s)|\le
 \frac{|s^{\overline{2p}}|C_{2p}}{(2p)!(\sigma+2p-1)}
                 N^{1-\sigma-2p}.}
\]
A rational upper bound for the finite rising-factorial product is used
for computed input boxes.  The remainder can be omitted to this error,
or computed more accurately by integrating finitely many unit cells and
using the same bound at a larger cutoff.

Thus for a particular $s\ne1$, choose $p$ with $\Re s>1-2p$ and a
certified separation from $s=1$, then increase $N$ until the error is small.
This computes a value even when the original series diverges.
Independence of $N$ follows from the finite summation identity.  Independence
of $p$ follows by two further integrations by parts in the remainder:
the extra correction is precisely the next displayed Bernoulli term,
and the far boundary vanishes in the common half-plane.  This proves
agreement without assuming an analytic continuation theorem.

These formulas are analytic there, with their only possible pole at $1$.
For the remainder, expand $t^{-w}$ near any fixed parameter.  The estimate
$|L(t)|^n/n!\le\rho^{-n}t^\rho$ gives the parameter-majorant proof
whenever $\rho<\Re s+2p-1$.  Hence the constructions also supply local
charts and all requested parameter derivatives.  The residue at $1$ is
$1$, read directly from $N^{1-s}/(s-1)$; its remaining part is computed
by the removable series quotient.

At $s=0$ take $p=1$; its rising factorials vanish and the formula gives
$\zeta(0)=-1/2$.  At $s=-1$ take $p=2$, so the remainder vanishes and
the finite terms give $\zeta(-1)=-1/12$.  This is a value of the proved
continuation computation, not a claim that the divergent sum of positive
integers converges to a negative number.  Standard zeta notation and
continuation formulas are recorded in
\href{https://dlmf.nist.gov/25.2}{DLMF, Section 25.2}.

\section{The common foundation is the finite comparison}
The constructions in this chapter use different kinds of error: an
exponential far tail, oscillatory integration by parts, Gaussian parameter
majorants, and a periodic-polynomial summation remainder.  None can be
replaced by a bare claim that the function is familiar.  Once its bound
has been proved, however, the existing interval, integration, derivative,
and equivalence constructions apply to it without rebuilding the theory.

\chapter{An Elliptic Integral and Its Differential Equation}
\label{ch:elliptic}

A final example ties together rational parametrization, algebraic-value
computations, integration, series, and a variable-coefficient fundamental
matrix.  The point is not to develop all elliptic functions, but to give
several equivalent ways of computing one useful family.

\section{Remove an endpoint singularity before computing}
For rational $0\le m<1$, the familiar expression is
\[
 K(m)=\int_0^1\frac{dx}{\sqrt{(1-x^2)(1-mx^2)}}.
\]
We do not begin by assuming this improper integral exists.  Use the rational
circle parameter $x=2u/(1+u^2)$.  For $0\le u<1$,
\[
 \sqrt{1-x^2}=\frac{1-u^2}{1+u^2},\qquad
 \frac{dx}{du}=\frac{2(1-u^2)}{(1+u^2)^2}.
\]
Cancellation suggests the bounded integral
\[
 \boxed{K(m)=\int_0^1
       \frac{2\,du}{\sqrt{(1+u^2)^2-4mu^2}}.}
\]
This is our first definition.  On $0\le m\le r<1$ its squared denominator
is at least $(1-r)(1+u^2)^2\ge1-r$.  Positive-root bisection and separated
reciprocals therefore give value computations, absolute bounds, and
continuity moduli for the integrand.

For truncated $u$-intervals the substitution theorem proves agreement with
the corresponding truncated original integral.  The omitted transformed
interval of length $\varepsilon$ contributes at most
$2\varepsilon/\sqrt{1-r}$, with any rational upper bound for this constant.
It follows that the original improper expression is an equivalent
computation, now accompanied by its missing tail estimate.

At $m=0$ the integrand is $2/(1+u^2)$, so
$K(0)=2A(1)=\pi/2$.

\section{A power series from the same integral}
Put $c_n=\binom{2n}{n}/4^n$, as in
Chapter~\ref{ch:infinite-series}.  The transformed integrand is
\[
 \frac2{1+u^2}\sum_{n=0}^{\infty}
       c_n m^n\left(\frac{2u}{1+u^2}\right)^{2n}.
\]
Since $0\le 2u/(1+u^2)\le1$ and $c_n\le1$, truncating after $N$ has
uniform error at most $2r^{N+1}/(1-r)$.  Integrating finite prefixes and
using that error bound is therefore justified.

Let
\[
 I_n=\int_0^1\frac2{1+u^2}
          \left(\frac{2u}{1+u^2}\right)^{2n}du.
\]
The angle substitution $x=2A(u)/\pi$ gives
$I_n=\pi\int_0^{1/2}S(x)^{2n}dx$.  Differentiate the finite product
$S^{2n-1}C$ using the already proved trigonometric derivatives and product
rule.  Its endpoint values vanish for $n\ge1$, so integration by parts
gives
\[
 2n I_n=(2n-1)I_{n-1},\qquad I_0=\pi/2.
\]
The powers and products have uniform finite increment estimates obtained
by multiplying those of $S$ and $C$, so the telescoping identity used here
is certified on the whole interval.  Hence $I_n=(\pi/2)c_n$ and
\[
 \boxed{K(m)=\frac\pi2\sum_{n=0}^{\infty}c_n^2m^n.}
\]
A tail bound is $(\pi/2)r^{N+1}/(1-r)$, which can be replaced by the
rational bound $2r^{N+1}/(1-r)$.

For example, at $m=1/2$ the first five terms of the normalized quantity
$2K(m)/\pi$ are
\[
 1+\frac18+\frac9{256}+\frac{25}{2048}
          +\frac{1225}{262144}.
\]
The remaining normalized tail is at most $1/16$.  More precise boxes
are obtained by the same rational recurrence, not by a new existence
argument.  Rectangle quadrature and this series now compute the same number.

\section{The differential equation is a coefficient identity}
Writing $a_n=(\pi/2)c_n^2$, the recurrence is
\[
 (n+1)^2a_{n+1}=(n+1/2)^2a_n.
\]
The differentiated tails converge uniformly on $m\le r<1$, by the power
series estimates.  Comparing coefficients therefore gives
\[
 \boxed{m(1-m)K''+(1-2m)K'-\tfrac14K=0,}
\]
with $K(0)=\pi/2$ and $K'(0)=\pi/8$.
This is proved from finite recurrences and their tails, not from an abstract
classification of period functions.

On an interval strictly inside $(0,1)$, put $Y=(K,K')^{\mathsf T}$.  It
satisfies the linear system
\[
 Y'=\begin{pmatrix}
 0&1\\
 \dfrac1{4m(1-m)}&\dfrac{2m-1}{m(1-m)}
 \end{pmatrix}Y.
\]
For example, on $[1/4,3/4]$ the row-sum norm of this coefficient matrix is
at most $4$.  The fundamental-matrix construction therefore has the
factorial majorant of the preceding chapter.  The series supplies
$Y(1/2)$, and propagation supplies $Y(3/4)$; integral-equation uniqueness
identifies it with direct series evaluation there.

The coefficient matrix is singular at $m=0$ and $m=1$.  We have not
applied its bounded-coefficient theorem across either point.  The series
handles the specified regular solution at $m=0$; each evolution computation
uses an explicitly separated interval.

\section{What the example connects}
We now have three actual computations of $K$: bounded quadrature of an
algebraic integrand, a rational coefficient series times $\pi/2$, and
fundamental-matrix propagation from a computed initial state.  Their
comparisons are proofs: rational substitution with an endpoint tail,
termwise finite integration with a uniform remainder, and uniqueness for
an integral equation.

The same matrix also computes a response to any finite list of impulses
in its two state coordinates on a separated interval.  Thus delta inputs
need no new theory for this example.  They are finite state updates
propagated by a computation already constructed for a concrete special
function.  This is the recurring pattern of the book: retain useful
representations, prove their equivalence, and reuse the resulting
computations in the next problem.

